\documentclass[11pt,a4paper]{article}

% ==================== TEMEL PAKETLER (pdflatex) ====================
\usepackage[T1]{fontenc}
\usepackage[utf8]{inputenc}
\usepackage{lmodern}
\usepackage[a4paper,margin=2cm,top=2.4cm,bottom=2.4cm]{geometry}
\usepackage{xcolor}
\usepackage{colortbl}
\usepackage{array}
\usepackage{booktabs}
\usepackage{longtable}
\usepackage{tabularx}
\usepackage{listings}
\usepackage{fancyhdr}
\usepackage{amsmath}
\usepackage{amssymb}
\usepackage{hyperref}

% ==================== RENK PALETİ ====================
\definecolor{anaMavi}{HTML}{132C4A}
\definecolor{vurguMavi}{HTML}{2E6DA4}
\definecolor{acikMavi}{HTML}{EAF2F9}
\definecolor{turuncu}{HTML}{D9701E}
\definecolor{yesil}{HTML}{2E7D32}
\definecolor{acikYesil}{HTML}{E8F5E9}
\definecolor{kirmizi}{HTML}{C0392B}
\definecolor{mor}{HTML}{6C3483}
\definecolor{acikMor}{HTML}{F3E9F7}
\definecolor{acikGri}{HTML}{F4F5F7}
\definecolor{ortaGri}{HTML}{6B7280}
\definecolor{koyuGri}{HTML}{2D2D2D}
\definecolor{kodArka}{HTML}{FBFBFA}

% ==================== HYPERREF ====================
\hypersetup{
  colorlinks=true,
  linkcolor={[HTML]{2E6DA4}},
  urlcolor={[HTML]{2E6DA4}},
  pdftitle={C++ ile Sayısal İşaret İşleme, Modülasyon ve Demodülasyon},
  pdfauthor={DSP Rehberi},
  bookmarksnumbered=true
}

% ==================== BAŞLIK STİLLERİ (titlesec olmadan) ====================
\setcounter{secnumdepth}{0}
\newcounter{secc}
\newcounter{subsecc}[secc]
\newcounter{subsubsecc}[subsecc]
\newcommand{\gsec}[1]{\par\phantomsection\stepcounter{secc}%
  \setcounter{subsecc}{0}%
  \vspace{18pt}{\color{anaMavi}\Large\bfseries\thesecc.\hspace{0.5em}#1}\par
  \vspace{2pt}{\color{turuncu}\rule{\linewidth}{1.6pt}}\par\vspace{8pt}%
  \addcontentsline{toc}{section}{\protect\numberline{\thesecc}#1}\markboth{#1}{#1}}
\newcommand{\gsub}[1]{\par\phantomsection\stepcounter{subsecc}%
  \setcounter{subsubsecc}{0}%
  \vspace{11pt}{\color{vurguMavi}\large\bfseries\thesecc.\thesubsecc\hspace{0.5em}#1}\par\vspace{4pt}%
  \addcontentsline{toc}{subsection}{\protect\numberline{\thesecc.\thesubsecc}#1}}
\newcommand{\gsubsub}[1]{\par\phantomsection\stepcounter{subsubsecc}%
  \vspace{7pt}{\color{koyuGri}\normalsize\bfseries #1}\par\vspace{2pt}}

% Kısım (Part) başlığı
\newcommand{\gpart}[2]{\clearpage
  \begin{center}\vspace*{1.5cm}
  {\color{turuncu}\rule{0.7\linewidth}{1.5pt}}\\[10pt]
  {\color{ortaGri}\Large KISIM #1}\\[8pt]
  {\color{anaMavi}\Huge\bfseries #2}\\[10pt]
  {\color{turuncu}\rule{0.7\linewidth}{1.5pt}}
  \end{center}\vspace{1cm}}

% ==================== BAŞLIK / ALTBİLGİ ====================
\setlength{\headheight}{14pt}
\pagestyle{fancy}
\fancyhf{}
\renewcommand{\headrulewidth}{0.4pt}
\renewcommand{\footrulewidth}{0.4pt}
\fancyhead[L]{\small\color{ortaGri}C++ ile Sayısal İşaret İşleme}
\fancyhead[R]{\small\color{ortaGri}\leftmark}
\fancyfoot[C]{\small\color{ortaGri}\thepage}

% ==================== KOD LİSTELEME (C++) ====================
\lstdefinestyle{cppstyle}{
  language=C++,
  basicstyle=\ttfamily\small\color{koyuGri},
  keywordstyle=\color{vurguMavi}\bfseries,
  commentstyle=\color{yesil}\itshape,
  stringstyle=\color{turuncu},
  numberstyle=\tiny\color{ortaGri},
  numbers=left,
  numbersep=8pt,
  backgroundcolor=\color{kodArka},
  frame=single,
  framesep=6pt,
  rulecolor=\color{ortaGri!40},
  breaklines=true,
  showstringspaces=false,
  tabsize=4,
  columns=fullflexible,
  keepspaces=true,
  extendedchars=true,
  literate={ı}{{\i}}1{İ}{{\.I}}1{ş}{{\c s}}1{Ş}{{\c S}}1{ğ}{{\u g}}1{Ğ}{{\u G}}1{ç}{{\c c}}1{Ç}{{\c C}}1{ö}{{\"o}}1{Ö}{{\"O}}1{ü}{{\"u}}1{Ü}{{\"U}}1{—}{{-{}-}}2{–}{{-}}1{’}{{'}}1{“}{{"}}1{”}{{"}}1,
  morekeywords={complex,vector,array,size_t,uint32_t,int32_t,uint16_t,int16_t,
    uint8_t,int8_t,uint64_t,int64_t,nullptr,constexpr,noexcept,override,final,
    auto,std,string,ifstream,ofstream,ios,M_PI,real,imag,norm,arg,conj,abs,
    Complex,Cplx,Signal,cf32,cf64}
}
\lstset{style=cppstyle}

% Kabuk/terminal stili
\lstdefinestyle{shstyle}{
  basicstyle=\ttfamily\small\color{koyuGri},
  backgroundcolor=\color{koyuGri!6},
  frame=single,framesep=6pt,rulecolor=\color{ortaGri!40},
  breaklines=true,showstringspaces=false,tabsize=4,columns=fullflexible,
  keepspaces=true,
  literate={ı}{{\i}}1{İ}{{\.I}}1{ş}{{\c s}}1{Ş}{{\c S}}1{ğ}{{\u g}}1{Ğ}{{\u G}}1{ç}{{\c c}}1{Ç}{{\c C}}1{ö}{{\"o}}1{Ö}{{\"O}}1{ü}{{\"u}}1{Ü}{{\"U}}1{—}{{-{}-}}2{–}{{-}}1,
}

% CMake stili
\lstdefinestyle{cmakestyle}{
  basicstyle=\ttfamily\small\color{koyuGri},
  keywordstyle=\color{mor}\bfseries,
  commentstyle=\color{yesil}\itshape,
  stringstyle=\color{turuncu},
  backgroundcolor=\color{acikMor!35},
  frame=single,framesep=6pt,rulecolor=\color{mor!40},
  breaklines=true,showstringspaces=false,tabsize=4,columns=fullflexible,
  keepspaces=true,
  morekeywords={cmake_minimum_required,project,add_executable,add_subdirectory,
    target_link_libraries,target_include_directories,find_package,set,
    add_library,target_compile_features,add_compile_options,foreach,endforeach,
    include,option,if,endif,PRIVATE,PUBLIC,VERSION,REQUIRED,LANGUAGES},
  literate={ı}{{\i}}1{İ}{{\.I}}1{ş}{{\c s}}1{Ş}{{\c S}}1{ğ}{{\u g}}1{Ğ}{{\u G}}1{ç}{{\c c}}1{Ç}{{\c C}}1{ö}{{\"o}}1{Ö}{{\"O}}1{ü}{{\"u}}1{Ü}{{\"U}}1,
}

% ==================== ÖZEL KUTULAR (tcolorbox olmadan) ====================
\newcommand{\callout}[4]{\par\smallskip\noindent\begingroup
  \setlength{\fboxsep}{8pt}\setlength{\fboxrule}{0.8pt}%
  \fcolorbox{#1}{#2}{\parbox{\dimexpr\linewidth-2\fboxsep-2\fboxrule\relax}{%
    {\bfseries\color{#1}#3}\par\smallskip #4}}\endgroup\par\smallskip}
\newcommand{\bilgi}[1]{\callout{vurguMavi}{acikMavi}{Bilgi}{#1}}
\newcommand{\ipucu}[1]{\callout{yesil}{acikYesil}{İpucu}{#1}}
\newcommand{\uyari}[1]{\callout{kirmizi}{kirmizi!7}{Dikkat}{#1}}
\newcommand{\ornek}[2][Örnek]{\callout{ortaGri!90!black}{acikGri}{#1}{#2}}
\newcommand{\notk}[2][Not]{\callout{mor}{acikMor}{#1}{#2}}

% Yardımcı komutlar
\newcommand{\code}[1]{\texttt{\color{koyuGri}#1}}
\renewcommand{\arraystretch}{1.35}

% ==================== BAŞLANGIÇ ====================
\begin{document}

% ---------- KAPAK ----------
\begin{titlepage}
\centering
\vspace*{1.4cm}
{\color{turuncu}\rule{\textwidth}{2pt}}\\[8pt]
{\color{anaMavi}\Huge\bfseries Sayısal İşaret İşleme}\\[8pt]
{\color{vurguMavi}\LARGE Modülasyon, Demodülasyon ve IQ İşleme}\\[8pt]
{\color{ortaGri}\large C++ ile Sıfırdan: DFT/FFT, Filtreler, AM/FM, PSK/QAM ve Tam Bir SDR Alıcısı}\\[6pt]
{\color{turuncu}\rule{\textwidth}{2pt}}\\[1.4cm]

\begin{center}
\fcolorbox{vurguMavi}{acikMavi}{\parbox{0.86\textwidth}{\centering\large
Bu belge önce \textbf{örnekleme, Nyquist ve karmaşık sayılar} gibi temelleri;
sonra \textbf{DFT/FFT ile frekans analizini}; ardından \textbf{FIR/IIR filtreleri};
sonra işin kalbi olan \textbf{IQ (karmaşık bant temeli)} kavramını; en sonda ise
\textbf{AM, FM, PM ve sayısal modülasyonların} tam modem örneklerini anlatır.
Kapanışta gerçek bir \textbf{FM radyo alıcısının} (RTL-SDR IQ akışından sese) uçtan
uca hattı yazılır. Tüm kod örnekleri İngilizce isimlendirmeyle, standart C++17 ile,
harici kütüphaneye ihtiyaç duymadan derlenecek şekilde yazılmıştır.\vspace{4pt}}}
\end{center}

\vspace{0.7cm}
\begin{center}
\fcolorbox{ortaGri!50}{acikGri}{\parbox{0.86\textwidth}{\small
\textbf{Kapsam --- Kısım 1 (Temeller):} işaret nedir \textbullet\ analog vs sayısal
\textbullet\ örnekleme \& Nyquist \textbullet\ aliasing \textbullet\ kuantalama \& ADC
\textbullet\ karmaşık sayılar \& fazörler \textbullet\ sinyal üretimi \& WAV G/Ç.\vspace{3pt}

\textbf{Kapsam --- Kısım 2 (Frekans Analizi):} Fourier fikri \textbullet\ DFT
\textbullet\ FFT (radix-2, özyinelemeli \& yinelemeli) \textbullet\ pencereleme
\textbullet\ genlik/faz spektrumu \textbullet\ güç spektral yoğunluğu.\vspace{3pt}

\textbf{Kapsam --- Kısım 3 (Filtreler):} konvolüsyon \textbullet\ FIR \& doğrusal faz
\textbullet\ pencere yöntemiyle tasarım \textbullet\ IIR \& biquad \textbullet\ kutup-sıfır
\textbullet\ hareketli ortalama \& CIC.\vspace{3pt}

\textbf{Kapsam --- Kısım 4 (IQ \& Bant Temeli):} analitik işaret \& Hilbert
\textbullet\ I/Q anlamı \textbullet\ NCO \& mikser ile frekans kaydırma
\textbullet\ desimasyon/interpolasyon \textbullet\ takım yıldızı (constellation).\vspace{3pt}

\textbf{Kapsam --- Kısım 5 (Modülasyon):} AM (zarf \& senkron) \textbullet\ FM (IQ diskriminatör)
\textbullet\ PM \textbullet\ ASK/FSK/PSK/QAM \textbullet\ QPSK modem \textbullet\ Costas \& PLL
\textbullet\ zamanlama kurtarma.\vspace{3pt}

\textbf{Kapsam --- Kısım 6 (Uçtan Uca Alıcı):} IQ dosya formatları \textbullet\ RTL-SDR akışı
\textbullet\ tam WBFM radyo hattı \textbullet\ stereo \& de-emphasis \textbullet\ performans notları.\vspace{3pt}

\textbf{Kapsam --- Kısım 7 (İleri):} stereo FM çözücü \textbullet\ AGC
\textbullet\ uçtan uca QPSK modem \textbullet\ OFDM (IFFT/FFT) \textbullet\ konvolüsyonel kod \& Viterbi.\vspace{3pt}

\textbf{Kapsam --- Kısım 8 (Projeler):} DTMF çözücü (Goertzel) \textbullet\ ses üzerinden FSK veri iletimi
\textbullet\ terminal spektrum analizörü \textbullet\ enstrüman akort cihazı \textbullet\ CMake ile derleme.\vspace{3pt}

\textbf{Kapsam --- Kısım 9 (Matematik):} diklik \& Euler \textbullet\ Fourier zinciri (FS/FT/DTFT/DFT)
\textbullet\ FFT ispatı \textbullet\ örnekleme teoreminin türetimi \textbullet\ konvolüsyon \& modülasyon
teoremleri \textbullet\ Parseval \textbullet\ Bessel \& Carson \textbullet\ analitik işaret \& anlık frekans.\vspace{3pt}}}
\end{center}

\vfill
{\color{ortaGri}\large Hazırlanma Tarihi: 20 Temmuz 2026}\\[4pt]
{\color{ortaGri} C++17 \textbullet\ \code{<complex>} \& \code{<vector>} \textbullet\ Nyquist--Shannon \textbullet\ RTL-SDR \textbullet\ CMake}
\vspace*{0.4cm}
\end{titlepage}

% ---------- İÇİNDEKİLER ----------
\tableofcontents
\thispagestyle{empty}
\clearpage

% #####################################################################
\gpart{1}{İşaret İşlemenin Temelleri}
% #####################################################################

% =====================================================================
\gsec{İşaret Nedir? Analogdan Sayısala}
% =====================================================================

Bir \textbf{işaret} (signal), bilgi taşıyan ve zamana (veya uzaya) göre değişen
bir büyüklüktür. Bir mikrofonun zara çarpan hava basıncını gerilime çevirmesi,
bir antenin havadaki elektromanyetik alanı akıma çevirmesi, bir sıcaklık
sensörünün ısıyı volta çevirmesi --- hepsi işaret üretir. İşaret işleme, bu
büyüklükleri \emph{ölçmek, dönüştürmek, filtrelemek, analiz etmek ve yeniden
üretmek} sanatıdır.

İşaretleri iki eksende sınıflandırırız:

\begin{center}
\begin{tabularx}{\textwidth}{@{}>{\bfseries\color{anaMavi}}l >{\raggedright\arraybackslash}X >{\raggedright\arraybackslash}X@{}}
\toprule
\rowcolor{acikMavi} & \color{anaMavi}\textbf{Zaman ekseni} & \color{anaMavi}\textbf{Genlik ekseni}\\
\midrule
\textbf{Analog} & sürekli ($t \in \mathbb{R}$) & sürekli (her değeri alabilir)\\
\textbf{Ayrık-zamanlı} & kesikli ($n \in \mathbb{Z}$) & sürekli\\
\textbf{Sayısal (dijital)} & kesikli & kesikli (kuantalanmış)\\
\bottomrule
\end{tabularx}
\end{center}

Gerçek dünya analogdur; bilgisayarlar ise sayısaldır. Aradaki köprü iki
bileşendir: \textbf{ADC} (Analog-to-Digital Converter) analog işareti belirli
aralıklarla ölçüp sayıya çevirir; \textbf{DAC} (Digital-to-Analog Converter) ise
sayı dizisini tekrar sürekli bir gerilime çevirir.

\bilgi{Sayısal işaret işlemenin (DSP) tüm gücü şu basit gerçekten gelir: Bir
işareti sayı dizisine çevirdiğimizde, onun üzerinde \emph{aritmetik} yapabiliriz.
Filtreleme bir çarpma-toplama, modülasyon bir çarpma, frekans analizi bir matris
dönüşümü hâline gelir. Analogda pahalı devrelerle yapılan her şey, sayısalda
birkaç satır koda dönüşür.}

\gsub{Ayrık-Zamanlı İşaretin Gösterimi}

Ayrık bir işareti $x[n]$ ile gösteririz; burada $n$ tamsayı örnek indeksidir.
Eğer örnekler $T_s$ saniye aralıkla alınıyorsa, \textbf{örnekleme frekansı}
$f_s = 1/T_s$'dir ve $n$'inci örnek gerçek zamanda $t = n\,T_s = n/f_s$ anına
karşılık gelir. C++'ta bir işaret çoğunlukla bir dizidir:

\begin{lstlisting}
#include <vector>
#include <cmath>

// Gerçek-değerli bir işaret genelde float/double vektörüdür.
using Signal = std::vector<float>;

// fs Hz örnekleme hızıyla, f Hz'lik saf bir sinüs üret.
// x[n] = A * sin(2*pi*f*n/fs + phi)
Signal makeSine(double freqHz, double fs, double durationSec,
                double amplitude = 1.0, double phaseRad = 0.0) {
    const size_t N = static_cast<size_t>(durationSec * fs);
    Signal x(N);
    const double w = 2.0 * M_PI * freqHz / fs;   // örnek başına faz artışı (rad)
    for (size_t n = 0; n < N; ++n)
        x[n] = static_cast<float>(amplitude * std::sin(w * n + phaseRad));
    return x;
}
\end{lstlisting}

\notk{Buradaki $w = 2\pi f/f_s$ niceliğine \emph{sayısal açısal frekans} denir
ve birimi ``radyan/örnek''tir. DSP'de neredeyse her formül $f_s$'e göre
normalize edilmiş bu $w$ cinsinden yazılır; çünkü algoritma açısından önemli olan
mutlak frekans değil, örnekleme hızına \emph{oranıdır}.}

% =====================================================================
\gsec{Örnekleme, Nyquist ve Aliasing}
% =====================================================================

DSP'nin en temel teoremi \textbf{Nyquist--Shannon örnekleme teoremidir}:
bant genişliği $B$ olan (yani en yüksek frekans bileşeni $B$ Hz olan) bir işareti
kayıpsız temsil etmek için örnekleme frekansı en az $2B$ olmalıdır:
\[
f_s > 2 f_{\max} \qquad(\text{Nyquist kriteri})
\]
$f_s/2$ değerine \textbf{Nyquist frekansı} denir ve sayısal dünyada temsil
edebileceğimiz en yüksek frekanstır.

\gsub{Aliasing: Kimlik Değiştiren Frekanslar}

Nyquist kriterini ihlal edersek --- yani $f_{\max}$'tan daha düşük bir hızla
örnekleme yaparsak --- yüksek frekanslar düşük frekanslar gibi \emph{kılık
değiştirir}. Buna \textbf{aliasing} (örtüşme) denir. Matematiksel olarak,
$f$ ve $f + k f_s$ frekansları (herhangi bir tamsayı $k$) örneklendiğinde
\emph{tamamen aynı} örnek dizisini üretir:
\[
\sin\!\big(2\pi (f + k f_s)\, n/f_s\big) = \sin\!\big(2\pi f\, n/f_s + 2\pi k n\big)
= \sin\!\big(2\pi f\, n/f_s\big)
\]
çünkü $2\pi k n$ terimi tam bir tur katıdır. Aşağıdaki kod bunu deneysel olarak
gösterir: $f_s = 1000$~Hz'te örneklenen 300~Hz ve 700~Hz'lik sinüsler bit-bit
aynı örnekleri verir (700 = 1000 $-$ 300, yani 700 Hz, $-300$ Hz'in aliasıdır).

\begin{lstlisting}
#include <cstdio>
#include <cmath>

int main() {
    const double fs = 1000.0;          // örnekleme frekansı
    const double f1 = 300.0, f2 = 700.0;
    printf(" n |   300 Hz   |   700 Hz   | fark\n");
    for (int n = 0; n < 8; ++n) {
        double a = std::sin(2*M_PI*f1*n/fs);
        double b = std::sin(2*M_PI*f2*n/fs);
        printf("%2d | %+.6f | %+.6f | %.1e\n", n, a, b, std::fabs(a-b));
    }
    // fark ~ 1e-16: iki farklı frekans, ayırt edilemez örnekler üretir!
    return 0;
}
\end{lstlisting}

\uyari{Aliasing geri döndürülemez bir bilgi kaybıdır. Örnekleme yapıldıktan
\emph{sonra} hangi frekansın ``gerçek'' olduğunu bilmenin yolu yoktur. Bu yüzden
her gerçek ADC'nin önünde, Nyquist üstü frekansları örneklemeden \emph{önce}
kesen bir analog \textbf{anti-aliasing (kıvrım önleyici) filtre} bulunur.}

\gsub{Aliasing'i Sese Çevirmek}

Aliasing'i somutlaştırmak için, $f_s/2$'yi aşan bir frekansın nereye
katlandığını hesaplayan küçük bir yardımcı yazalım. Katlanan frekans
$f_{\text{alias}} = |f - k f_s|$ olup, sonucu $[0, f_s/2]$ aralığına indirir:

\begin{lstlisting}
#include <cmath>

// Bir f frekansının fs ile örneklendiğinde nereye "alias" olacağını verir.
double aliasedFrequency(double f, double fs) {
    double folded = std::fmod(f, fs);           // [0, fs) aralığına indir
    if (folded < 0) folded += fs;
    if (folded > fs / 2.0) folded = fs - folded; // Nyquist'in üstünü katla
    return folded;
}

// Örnek: fs = 48 kHz
//   aliasedFrequency(1000,  48000) -> 1000  (sorun yok)
//   aliasedFrequency(47000, 48000) -> 1000  (47 kHz, 1 kHz gibi duyulur!)
//   aliasedFrequency(25000, 48000) -> 23000 (Nyquist'in hemen üstü katlanır)
\end{lstlisting}

% =====================================================================
\gsec{Kuantalama ve ADC}
% =====================================================================

Örnekleme zaman eksenini kesikli yapar; \textbf{kuantalama} ise genlik eksenini
kesikli yapar. Bir $B$-bitlik ADC, giriş aralığını $2^B$ seviyeye böler. Örneğin
16-bit bir ADC, gerilim aralığını $65536$ basamağa böler. Her örnek en yakın
basamağa yuvarlanır; aradaki fark \textbf{kuantalama gürültüsüdür}.

Bir ideal kuantalayıcının sinyal-gürültü oranı (SNR) yaklaşık şudur:
\[
\text{SNR}_{\text{dB}} \approx 6{,}02\,B + 1{,}76
\]
Yani her ek bit yaklaşık 6~dB dinamik aralık kazandırır. 16-bit $\approx 98$~dB,
12-bit (tipik SDR) $\approx 74$~dB.

\begin{lstlisting}
#include <cmath>
#include <cstdint>

// [-1, 1] aralığındaki float'ı B-bit işaretli tamsayıya kuantala.
int32_t quantize(float x, int bits) {
    const int32_t levels = (1 << (bits - 1)) - 1;   // örn. 16-bit -> 32767
    float clamped = std::fmax(-1.0f, std::fmin(1.0f, x));
    return static_cast<int32_t>(std::lround(clamped * levels));
}

// Kuantalanmış değeri geri float'a çevir (DAC'ın yaptığı).
float dequantize(int32_t q, int bits) {
    const int32_t levels = (1 << (bits - 1)) - 1;
    return static_cast<float>(q) / levels;
}
\end{lstlisting}

\ipucu{RTL-SDR gibi ucuz alıcılar 8-bit örnekler verir (her I ve Q için bir
bayt, \code{0..255}, orta nokta 127.5). Bu yüzden bir SDR programının ilk adımı
neredeyse her zaman bu baytları $[-1, 1]$ aralığına ölçekleyip \code{float}'a
çevirmektir. Bunun kodunu Kısım 6'da göreceğiz.}

% =====================================================================
\gsec{Karmaşık Sayılar ve Fazörler}
% =====================================================================

Modülasyon ve IQ işlemenin dilini konuşabilmek için \textbf{karmaşık sayılar}
şarttır. Bir karmaşık sayı $z = a + jb$'dir; burada $j = \sqrt{-1}$ (mühendisler
$i$ yerine $j$ kullanır, çünkü $i$ akımı temsil eder). Gerçek kısım $a =
\Re\{z\}$, sanal kısım $b = \Im\{z\}$'dir.

En kritik ilişki \textbf{Euler formülüdür}:
\[
e^{j\theta} = \cos\theta + j\sin\theta
\]
Bu, birim çember üzerinde $\theta$ açısındaki bir noktadır. Zamanla dönen bir
karmaşık üstel $e^{j 2\pi f t}$'ye \textbf{fazör} denir ve $f$ frekansında dönen
bir vektör olarak düşünülür. İşte modülasyonun tüm geometrisi budur: bilgiyi bu
dönen vektörün \emph{genliğine}, \emph{açısına} veya \emph{dönme hızına}
gömeriz.

\begin{center}
\begin{tabularx}{\textwidth}{@{}>{\bfseries\color{anaMavi}}l >{\raggedright\arraybackslash}X@{}}
\toprule
\rowcolor{acikMavi}\color{anaMavi}\textbf{Büyüklük} & \color{anaMavi}\textbf{Anlamı ve C++ karşılığı}\\
\midrule
Genlik (magnitude) & $|z| = \sqrt{a^2+b^2}$ --- \code{std::abs(z)}. AM'de bilgi burada.\\
Faz (argument) & $\angle z = \operatorname{atan2}(b,a)$ --- \code{std::arg(z)}. PM/PSK'de bilgi burada.\\
Anlık frekans & fazın türevi $\frac{d}{dt}\angle z$. FM'de bilgi burada.\\
Eşlenik & $z^* = a - jb$ --- \code{std::conj(z)}. Frekans kaydırmada kullanılır.\\
Güç & $|z|^2 = a^2+b^2$ --- \code{std::norm(z)} (dikkat: karekök \emph{yok}).\\
\bottomrule
\end{tabularx}
\end{center}

\uyari{C++ standart kütüphanesinde \code{std::norm(z)} matematiksel normu
(uzunluğu) \emph{değil}, uzunluğun \emph{karesini} ($|z|^2$) döndürür. Uzunluk
için \code{std::abs(z)} kullanın. Bu, DSP kodunda sık yapılan bir hatadır.}

\gsub{C++'ta \code{std::complex} ile Çalışmak}

Standart \code{<complex>} başlığı bize hazır bir karmaşık aritmetik verir. SDR
dünyasında yaygın olarak \code{std::complex<float>} kullanılır (bellek ve hız
için); GNU Radio'daki \code{gr\_complex} de tam olarak budur.

\begin{lstlisting}
#include <complex>
#include <cstdio>

using cf32 = std::complex<float>;   // I = real, Q = imag

int main() {
    cf32 z(3.0f, 4.0f);             // 3 + 4j  (I=3, Q=4)

    printf("z        = %.1f + %.1fj\n", z.real(), z.imag());
    printf("|z|      = %.4f\n", std::abs(z));    // 5.0 (genlik)
    printf("|z|^2    = %.4f\n", std::norm(z));   // 25.0 (guc, karekoksuz!)
    printf("faz      = %.4f rad\n", std::arg(z)); // atan2(4,3) = 0.9273
    printf("eslenik  = %.1f + %.1fj\n",
           std::conj(z).real(), std::conj(z).imag()); // 3 - 4j

    // Euler: bir birim fazor uret (45 derece)
    cf32 p = std::polar(1.0f, 3.14159265f / 4.0f);  // e^(j*pi/4)
    printf("polar    = %.4f + %.4fj\n", p.real(), p.imag()); // ~0.707 + 0.707j

    // Iki fazoru carpmak aci'larini TOPLAR (frekans kaydirmanin ozu)
    cf32 a = std::polar(1.0f, 0.3f);
    cf32 b = std::polar(1.0f, 0.5f);
    printf("carpim aci = %.4f (=0.8?)\n", std::arg(a * b));  // 0.8
    return 0;
}
\end{lstlisting}

\bilgi{Neden karmaşık sayılar? Çünkü \textbf{çarpım açıları toplar}: $e^{j\alpha}
\cdot e^{j\beta} = e^{j(\alpha+\beta)}$. Bir işareti $e^{-j 2\pi f_0 t/f_s}$ ile
çarpmak, onu frekans ekseninde $-f_0$ kadar kaydırır --- işte tüm mikserlerin,
kanal seçicilerin ve demodülatörlerin çalışma prensibi budur. Gerçek sinüslerle
bunu yapmaya çalışsaydık, çarpım hem toplam hem fark frekansını üretir ve
istenmeyen bileşeni filtrelememiz gerekirdi.}

% =====================================================================
\gsec{Sinyal Üretimi ve WAV Dosyaları}
% =====================================================================

Sonuçlarımızı duyabilmek ve grafiklemek için sinyalleri diske yazmak
gerekir. En yaygın format \textbf{WAV}'dır (RIFF konteyneri, PCM veri).
Aşağıda harici kütüphane olmadan 16-bit mono/stereo WAV yazan tam bir yardımcı
var --- bu belge boyunca üretilen sesleri kaydetmek için kullanacağız.

\begin{lstlisting}
#include <cstdint>
#include <cstdio>
#include <fstream>
#include <vector>
#include <cmath>

// 16-bit PCM WAV yazici. samples: [-1,1] araligindaki interleaved kanallar.
void writeWav(const std::string& path, const std::vector<float>& samples,
              uint32_t sampleRate, uint16_t channels = 1) {
    std::ofstream f(path, std::ios::binary);
    const uint32_t nSamples = static_cast<uint32_t>(samples.size());
    const uint16_t bitsPerSample = 16;
    const uint16_t blockAlign = channels * bitsPerSample / 8;
    const uint32_t byteRate = sampleRate * blockAlign;
    const uint32_t dataBytes = nSamples * sizeof(int16_t);

    auto put32 = [&](uint32_t v){ f.write(reinterpret_cast<char*>(&v), 4); };
    auto put16 = [&](uint16_t v){ f.write(reinterpret_cast<char*>(&v), 2); };

    f.write("RIFF", 4);  put32(36 + dataBytes);  f.write("WAVE", 4);
    f.write("fmt ", 4);  put32(16);              // fmt yigin boyutu
    put16(1);                                     // PCM
    put16(channels);     put32(sampleRate);       put32(byteRate);
    put16(blockAlign);   put16(bitsPerSample);
    f.write("data", 4);  put32(dataBytes);

    for (float s : samples) {
        float c = std::fmax(-1.0f, std::fmin(1.0f, s));   // kirp (clip)
        int16_t q = static_cast<int16_t>(std::lround(c * 32767.0f));
        f.write(reinterpret_cast<char*>(&q), sizeof(q));
    }
}

int main() {
    const double fs = 48000.0;
    std::vector<float> tone;
    // 2 saniye, 440 Hz (La notasi), yavasca sonlanan zarfla
    for (int n = 0; n < static_cast<int>(2 * fs); ++n) {
        double t = n / fs;
        double env = std::exp(-1.5 * t);                 // ussel sonumlenme
        tone.push_back(static_cast<float>(0.6 * env * std::sin(2*M_PI*440*t)));
    }
    writeWav("la440.wav", tone, 48000, 1);
    return 0;
}
\end{lstlisting}

\ipucu{Ürettiğiniz herhangi bir işaretin frekans içeriğini hızlıca görmek için
onu WAV'a yazıp Audacity veya \code{sox} ile spektrogramına bakabilirsiniz:
\code{sox sinyal.wav -n spectrogram}. Bir demodülatörü hata ayıklarken
``kulakla dinlemek'' şaşırtıcı derecede güçlü bir araçtır.}

% #####################################################################
\gpart{2}{Frekans Alanı: DFT ve FFT}
% #####################################################################

% =====================================================================
\gsec{Fourier Fikri}
% =====================================================================

İşaret işlemenin en derin fikri şudur: \textbf{her işaret, farklı frekanslardaki
sinüslerin toplamı olarak yazılabilir}. Zaman alanında karmaşık görünen bir
sinyal, frekans alanında birkaç basit tepe hâline gelebilir. Filtreleme,
modülasyon, sıkıştırma --- çoğu işlem frekans alanında düşünüldüğünde apaçık
olur.

\textbf{Ayrık Fourier Dönüşümü (DFT)}, $N$ örneklik bir işareti $N$ frekans
bileşenine dönüştürür:
\[
X[k] = \sum_{n=0}^{N-1} x[n]\, e^{-j 2\pi k n / N}, \qquad k = 0,1,\dots,N-1
\]
Buradaki $X[k]$ karmaşıktır: büyüklüğü $|X[k]|$ o frekanstaki bileşenin gücünü,
açısı $\angle X[k]$ ise fazını verir. $k$'ıncı bileşenin karşılık geldiği gerçek
frekans $f_k = k f_s / N$'dir.

\gsub{DFT'nin Doğrudan Uygulaması}

Tanımı olduğu gibi koda dökmek $O(N^2)$'lik bir algoritma verir --- yavaştır ama
her zaman doğrudur ve FFT'yi doğrulamak için harika bir referanstır:

\begin{lstlisting}
#include <complex>
#include <vector>
#include <cmath>

using cf64 = std::complex<double>;

// Dogrudan DFT: O(N^2). Girdi gercek veya karmasik olabilir.
std::vector<cf64> dft(const std::vector<cf64>& x) {
    const size_t N = x.size();
    std::vector<cf64> X(N);
    for (size_t k = 0; k < N; ++k) {
        cf64 sum(0.0, 0.0);
        for (size_t n = 0; n < N; ++n) {
            double angle = -2.0 * M_PI * k * n / N;
            sum += x[n] * cf64(std::cos(angle), std::sin(angle));
        }
        X[k] = sum;
    }
    return X;
}

// Ters DFT (IDFT): frekanstan zamana geri don.
std::vector<cf64> idft(const std::vector<cf64>& X) {
    const size_t N = X.size();
    std::vector<cf64> x(N);
    for (size_t n = 0; n < N; ++n) {
        cf64 sum(0.0, 0.0);
        for (size_t k = 0; k < N; ++k) {
            double angle = +2.0 * M_PI * k * n / N;   // isaret ters, +
            sum += X[k] * cf64(std::cos(angle), std::sin(angle));
        }
        x[n] = sum / static_cast<double>(N);          // 1/N normalizasyonu
    }
    return x;
}
\end{lstlisting}

% =====================================================================
\gsec{FFT: Hızlı Fourier Dönüşümü}
% =====================================================================

DFT'nin $O(N^2)$ maliyeti büyük $N$ için dayanılmaz olur ($N = 10^6$ için
$10^{12}$ işlem). \textbf{FFT}, Cooley--Tukey algoritmasıyla bunu
$O(N \log N)$'e düşürür --- $N = 10^6$ için milyon kat hızlanma. Fikir
\emph{böl ve fethet}: $N$ noktalı DFT'yi çift ve tek indeksli örneklerin iki
adet $N/2$ noktalı DFT'sine ayır, sonra ``kelebek'' (butterfly) işlemleriyle
birleştir.

\gsub{Özyinelemeli Radix-2 FFT}

En anlaşılır uygulama özyinelemelidir. $N$'in 2'nin kuvveti olması gerekir:

\begin{lstlisting}
#include <complex>
#include <vector>
#include <cmath>

using cf64 = std::complex<double>;

// Ozyinelemeli radix-2 DIT FFT. x.size() 2'nin kuvveti olmali.
void fftRecursive(std::vector<cf64>& x) {
    const size_t N = x.size();
    if (N <= 1) return;

    // Cift ve tek indeksleri ayir
    std::vector<cf64> even(N/2), odd(N/2);
    for (size_t i = 0; i < N/2; ++i) {
        even[i] = x[2*i];
        odd[i]  = x[2*i + 1];
    }
    fftRecursive(even);
    fftRecursive(odd);

    // Kelebek birlestirme
    for (size_t k = 0; k < N/2; ++k) {
        double angle = -2.0 * M_PI * k / N;
        cf64 tw = std::polar(1.0, angle) * odd[k];   // twiddle carpani
        x[k]        = even[k] + tw;
        x[k + N/2]  = even[k] - tw;
    }
}
\end{lstlisting}

\gsub{Yerinde (In-Place) Yinelemeli FFT}

Üretim kodu özyineleme yerine, bit-ters (bit-reversal) permütasyon + iç içe
döngü kullanır. Bellek tahsisi yapmaz, önbellek dostudur ve çok daha hızlıdır:

\begin{lstlisting}
#include <complex>
#include <vector>
#include <cmath>

using cf64 = std::complex<double>;

// Yinelemeli, yerinde radix-2 FFT. invert=true ise ters FFT (1/N'siz).
void fft(std::vector<cf64>& a, bool invert = false) {
    const size_t n = a.size();
    if (n <= 1) return;

    // 1) Bit-ters permutasyon
    for (size_t i = 1, j = 0; i < n; ++i) {
        size_t bit = n >> 1;
        for (; j & bit; bit >>= 1) j ^= bit;
        j ^= bit;
        if (i < j) std::swap(a[i], a[j]);
    }

    // 2) Kelebekleri katman katman uygula (len = 2, 4, 8, ...)
    for (size_t len = 2; len <= n; len <<= 1) {
        double ang = (invert ? +2.0 : -2.0) * M_PI / len;
        cf64 wlen(std::cos(ang), std::sin(ang));
        for (size_t i = 0; i < n; i += len) {
            cf64 w(1.0, 0.0);
            for (size_t k = 0; k < len/2; ++k) {
                cf64 u = a[i + k];
                cf64 v = a[i + k + len/2] * w;
                a[i + k]         = u + v;
                a[i + k + len/2] = u - v;
                w *= wlen;
            }
        }
    }

    if (invert)
        for (cf64& x : a) x /= static_cast<double>(n);  // ters icin 1/N
}
\end{lstlisting}

\ornek[Kullanım: Bir Frekansı Tespit Et]{Aşağıdaki program, 8000~Hz'te
örneklenmiş ve içinde 1000~Hz'lik bir ton bulunan bir işaretin FFT'sini alıp en
güçlü frekans binini bulur.}

\begin{lstlisting}
#include <cstdio>

int main() {
    const double fs = 8000.0;
    const size_t N = 1024;                       // 2'nin kuvveti
    std::vector<cf64> x(N);
    for (size_t n = 0; n < N; ++n)
        x[n] = std::sin(2*M_PI*1000.0*n/fs);     // gercek girdi, Q=0

    fft(x);                                       // yerinde donusum

    // En buyuk buyuklukteki bin'i bul (yalniz ilk yariya bak: gercek girdi)
    size_t peak = 0; double best = 0;
    for (size_t k = 0; k < N/2; ++k) {
        double mag = std::abs(x[k]);
        if (mag > best) { best = mag; peak = k; }
    }
    printf("Tepe bin = %zu -> %.1f Hz\n", peak, peak * fs / N); // ~1000 Hz
    return 0;
}
\end{lstlisting}

\notk[Frekans çözünürlüğü]{FFT'nin bin aralığı $\Delta f = f_s / N$'dir. Daha ince
frekans çözünürlüğü için ya daha uzun kayıt alırsınız (büyük $N$) ya da işareti
sıfırlarla doldurursunuz (zero-padding). Bin çözünürlüğü zamanla ters
orantılıdır: 1~Hz çözünürlük için en az 1 saniyelik kayıt gerekir.}

% =====================================================================
\gsec{Pencereleme (Windowing)}
% =====================================================================

DFT, işaretin \emph{periyodik} olduğunu varsayar. Ama biz sonlu bir parça
kesip aldığımızda, kesme noktalarındaki süreksizlik frekans alanında
\textbf{spektral sızıntıya} (spectral leakage) yol açar: keskin bir tepe olması
gereken yerde, yanlara yayılan yan loblar (sidelobes) oluşur. Çözüm, işareti
kenarlarda yumuşakça sıfıra indiren bir \textbf{pencere fonksiyonuyla} çarpmaktır.

\begin{center}
\begin{tabularx}{\textwidth}{@{}>{\bfseries\color{anaMavi}}l >{\raggedright\arraybackslash}X@{}}
\toprule
\rowcolor{acikMavi}\color{anaMavi}\textbf{Pencere} & \color{anaMavi}\textbf{Özelliği}\\
\midrule
Dikdörtgen (yok) & En dar ana lob, en kötü sızıntı. Geçici olayları ölçmek için.\\
Hann & İyi genel amaçlı denge; yumuşak, $-31$~dB yan lob.\\
Hamming & Hann'a benzer, ilk yan lobu bastırır ($-43$~dB).\\
Blackman & Çok düşük yan lob ($-58$~dB), daha geniş ana lob. Zayıf tonları görmek için.\\
\bottomrule
\end{tabularx}
\end{center}

\begin{lstlisting}
#include <vector>
#include <cmath>

// Hann penceresi: w[n] = 0.5 * (1 - cos(2*pi*n/(N-1)))
std::vector<double> hannWindow(size_t N) {
    std::vector<double> w(N);
    for (size_t n = 0; n < N; ++n)
        w[n] = 0.5 * (1.0 - std::cos(2.0 * M_PI * n / (N - 1)));
    return w;
}

// Blackman penceresi: daha iyi yan-lob bastirma
std::vector<double> blackmanWindow(size_t N) {
    std::vector<double> w(N);
    for (size_t n = 0; n < N; ++n) {
        double t = 2.0 * M_PI * n / (N - 1);
        w[n] = 0.42 - 0.5 * std::cos(t) + 0.08 * std::cos(2 * t);
    }
    return w;
}

// Spektrum almadan once pencereyi uygula.
void applyWindow(std::vector<cf64>& x, const std::vector<double>& w) {
    for (size_t n = 0; n < x.size(); ++n) x[n] *= w[n];
}
\end{lstlisting}

\gsub{Güç Spektral Yoğunluğu (PSD)}

Bir işaretin gücünün frekansa nasıl dağıldığını görmek için PSD hesaplarız.
Pratikte pencerelenmiş FFT'nin büyüklük karesini alıp desibele çeviririz. SDR
programlarındaki ``şelale'' (waterfall) görüntüsünün altında yatan hesap budur:

\begin{lstlisting}
#include <vector>
#include <cmath>

// Tek bir kare icin dB cinsinden guc spektrumu.
std::vector<double> powerSpectrumDb(std::vector<cf64> x,
                                    const std::vector<double>& win) {
    const size_t N = x.size();
    applyWindow(x, win);
    fft(x);                              // (onceki bolumdeki FFT)

    // Pencere kazancini telafi et
    double winSum = 0; for (double w : win) winSum += w;

    std::vector<double> psd(N);
    for (size_t k = 0; k < N; ++k) {
        double mag = std::abs(x[k]) / winSum;
        psd[k] = 20.0 * std::log10(mag + 1e-12);  // +eps: log(0)'dan kacin
    }
    return psd;
}
\end{lstlisting}

% #####################################################################
\gpart{3}{Filtreler: FIR ve IIR}
% #####################################################################

% =====================================================================
\gsec{Konvolüsyon ve FIR Filtreler}
% =====================================================================

Bir \textbf{filtre}, işaretin belli frekanslarını geçirip diğerlerini
zayıflatan bir sistemdir. En temel yapı taşı \textbf{konvolüsyondur}: çıktı,
girdinin filtre \emph{dürtü yanıtı} $h[n]$ ile konvolüsyonudur:
\[
y[n] = \sum_{k=0}^{M-1} h[k]\, x[n-k]
\]
Dürtü yanıtı sonlu uzunluktaysa (yani $h$ sonlu bir katsayı dizisiyse) buna
\textbf{FIR} (Finite Impulse Response) filtre denir. FIR filtreler her zaman
kararlıdır ve \textbf{tam doğrusal faz} sağlayabilir (tüm frekansları aynı
gecikmeyle geciktirir; ses ve haberleşmede kritik).

\gsub{Temel FIR Uygulaması}

\begin{lstlisting}
#include <vector>

// Zamanla akan bir FIR filtresi. Katsayilar (taps) sabit; her ornek icin
// gecmis M orneklik pencereyi bir dairesel tampon (ring buffer) ile tutar.
class FirFilter {
    std::vector<float> taps_;    // h[k] katsayilari
    std::vector<float> buf_;     // gecmis ornekler (dairesel)
    size_t pos_ = 0;
public:
    explicit FirFilter(std::vector<float> taps)
        : taps_(std::move(taps)), buf_(taps_.size(), 0.0f) {}

    // Bir ornek isle, bir ornek dondur.
    float process(float x) {
        buf_[pos_] = x;
        float acc = 0.0f;
        size_t idx = pos_;
        for (size_t k = 0; k < taps_.size(); ++k) {
            acc += taps_[k] * buf_[idx];
            idx = (idx == 0) ? taps_.size() - 1 : idx - 1;  // geriye sar
        }
        pos_ = (pos_ + 1) % taps_.size();
        return acc;
    }

    void reset() { std::fill(buf_.begin(), buf_.end(), 0.0f); pos_ = 0; }
};
\end{lstlisting}

\gsub{Pencere Yöntemiyle Alçak Geçiren FIR Tasarımı}

Peki katsayıları nereden buluruz? İdeal bir alçak geçiren filtrenin dürtü yanıtı
bir \textbf{sinc} fonksiyonudur. Onu sonlu uzunlukta kesip bir pencereyle
yumuşatarak pratik bir filtre elde ederiz:

\begin{lstlisting}
#include <vector>
#include <cmath>

// Pencere-sinc yontemiyle alcak geciren FIR tasarla.
//   cutoffHz : -6 dB kesim frekansi
//   fs       : ornekleme frekansi
//   numTaps  : katsayi sayisi (tek sayi onerilir -> simetrik, tam dogrusal faz)
std::vector<float> designLowpass(double cutoffHz, double fs, size_t numTaps) {
    std::vector<float> h(numTaps);
    const double fc = cutoffHz / fs;          // normalize kesim (0..0.5)
    const double mid = (numTaps - 1) / 2.0;
    double sum = 0.0;

    for (size_t n = 0; n < numTaps; ++n) {
        double m = n - mid;
        // Ideal alcak geciren dPurtu yaniti: sinc
        double sinc = (m == 0.0) ? 2.0 * fc
                                 : std::sin(2.0 * M_PI * fc * m) / (M_PI * m);
        // Hamming penceresi
        double win = 0.54 - 0.46 * std::cos(2.0 * M_PI * n / (numTaps - 1));
        h[n] = static_cast<float>(sinc * win);
        sum += h[n];
    }
    // DC'de birim kazanc icin normalize et
    for (float& c : h) c /= static_cast<float>(sum);
    return h;
}

// Yuksek geciren = all-pass eksi alcak geciren (spektral tersleme)
std::vector<float> designHighpass(double cutoffHz, double fs, size_t numTaps) {
    auto h = designLowpass(cutoffHz, fs, numTaps);
    for (float& c : h) c = -c;                 // isaretini ters cevir
    h[(numTaps - 1) / 2] += 1.0f;              // merkeze birim dPurtu ekle
    return h;
}
\end{lstlisting}

\ornek[Uygulama: Gürültülü Sinyali Temizle]{500~Hz'lik yararlı tonu 5000~Hz'lik
gürültüden ayıralım.}

\begin{lstlisting}
int main() {
    const double fs = 44100.0;
    // 500 Hz sinyal + 5000 Hz gurultu
    std::vector<float> in;
    for (int n = 0; n < 4410; ++n) {
        double t = n / fs;
        in.push_back(0.5*std::sin(2*M_PI*500*t) + 0.5*std::sin(2*M_PI*5000*t));
    }
    // 1500 Hz'te kesen 101-tap alcak geciren filtre
    FirFilter lp(designLowpass(1500.0, fs, 101));
    std::vector<float> out;
    for (float x : in) out.push_back(lp.process(x));
    // out: yalnizca 500 Hz kalir; 5000 Hz ~40 dB bastirilmis olur.
    return 0;
}
\end{lstlisting}

% =====================================================================
\gsec{IIR Filtreler ve Biquad}
% =====================================================================

\textbf{IIR} (Infinite Impulse Response) filtreler geri besleme kullanır: çıktı
hem geçmiş girdilere hem geçmiş \emph{çıktılara} bağlıdır. Bu sayede çok az
katsayıyla çok keskin filtreler kurulabilir --- ama kararlılık artık garanti
değildir ve faz doğrusal olmaz. En yaygın yapı taşı \textbf{biquad}'dır (ikinci
dereceden bölüm):
\[
y[n] = b_0 x[n] + b_1 x[n\!-\!1] + b_2 x[n\!-\!2]
       - a_1 y[n\!-\!1] - a_2 y[n\!-\!2]
\]

\begin{lstlisting}
#include <cmath>

// Transpoze Direct Form II biquad (sayisal olarak daha kararli).
struct Biquad {
    float b0=1, b1=0, b2=0, a1=0, a2=0;   // katsayilar (a0=1'e normalize)
    float z1=0, z2=0;                     // durum degiskenleri

    float process(float x) {
        float y = b0 * x + z1;
        z1 = b1 * x - a1 * y + z2;
        z2 = b2 * x - a2 * y;
        return y;
    }
    void reset() { z1 = z2 = 0; }

    // RBJ "Audio EQ Cookbook" katsayilariyla alcak geciren biquad kur.
    static Biquad lowpass(double freqHz, double fs, double Q = 0.7071) {
        Biquad bq;
        double w0 = 2.0 * M_PI * freqHz / fs;
        double cw = std::cos(w0), sw = std::sin(w0);
        double alpha = sw / (2.0 * Q);
        double a0 = 1 + alpha;
        bq.b0 = (1 - cw) / 2 / a0;
        bq.b1 = (1 - cw)     / a0;
        bq.b2 = (1 - cw) / 2 / a0;
        bq.a1 = (-2 * cw)    / a0;
        bq.a2 = (1 - alpha)  / a0;
        return bq;
    }
    // Band geciren (tepe kazancli) biquad
    static Biquad bandpass(double freqHz, double fs, double Q = 5.0) {
        Biquad bq;
        double w0 = 2.0 * M_PI * freqHz / fs;
        double cw = std::cos(w0), sw = std::sin(w0);
        double alpha = sw / (2.0 * Q);
        double a0 = 1 + alpha;
        bq.b0 =  alpha / a0;  bq.b1 = 0;  bq.b2 = -alpha / a0;
        bq.a1 = (-2 * cw) / a0;  bq.a2 = (1 - alpha) / a0;
        return bq;
    }
};
\end{lstlisting}

\uyari{IIR filtreler tasarımı kolay ama tuzaklıdır. Katsayılar kutupları
(poles) birim çemberin \emph{dışına} taşırsa filtre patlar (kararsız olur).
Yüksek dereceli IIR'ları tek bir denklemle değil, art arda bağlı biquad'larla
(cascade) uygulayın --- bu, sayısal hata birikimini ve kararsızlığı ciddi
biçimde azaltır.}

\gsub{Hareketli Ortalama ve CIC: Ucuz Filtreler}

En ucuz alçak geçiren filtre \textbf{hareketli ortalamadır}. Onu bir toplam-tut
(running sum) ile $O(1)$'de uygulayabiliriz --- SDR'da desimasyon öncesi kaba
filtreleme için idealdir:

\begin{lstlisting}
#include <vector>

// O(1) hareketli ortalama filtresi (kutu filtre / boxcar).
class MovingAverage {
    std::vector<float> buf_;
    size_t pos_ = 0, count_ = 0;
    float sum_ = 0.0f;
public:
    explicit MovingAverage(size_t window) : buf_(window, 0.0f) {}
    float process(float x) {
        sum_ -= buf_[pos_];              // pencereden cikani cikar
        buf_[pos_] = x;
        sum_ += x;                       // yeni gireni ekle
        pos_ = (pos_ + 1) % buf_.size();
        if (count_ < buf_.size()) ++count_;
        return sum_ / count_;
    }
};
\end{lstlisting}

% #####################################################################
\gpart{4}{IQ ve Karmaşık Bant Temeli}
% #####################################################################

% =====================================================================
\gsec{IQ Nedir? Analitik İşaret}
% =====================================================================

Şimdi tüm modern haberleşmenin kalbine geliyoruz. \textbf{IQ}, ``In-phase''
(I, eş-fazlı) ve ``Quadrature'' (Q, dik-fazlı, $90^\circ$ kaymış) bileşenlerin kısa
adıdır. Bir IQ örneği aslında bir \textbf{karmaşık sayıdır}: $z = I + jQ$.

Neden gerçek bir sinüs yerine karmaşık bir sayı? Çünkü gerçek bir sinüs
$\cos(2\pi f t)$, frekans alanında hem $+f$ hem $-f$'te iki tepe içerir --- yani
pozitif ve negatif frekansı ayırt edemez. Oysa karmaşık üstel $e^{j 2\pi f t}$
yalnızca $+f$'te tek bir tepe içerir. Bu ``tek taraflılık'' sayesinde:

\begin{itemize}
\item Bir sinyalin frekansının \emph{işaretini} (yönünü) bilebiliriz.
\item Frekansı yukarı/aşağı kaydırırken görüntü (image) frekansı oluşmaz.
\item Genlik, faz ve frekansı aynı anda, ayrı ayrı çekebiliriz.
\end{itemize}

\bilgi{Bir SDR alıcısı antenden gelen yüksek frekanslı gerçek işareti, donanımda
iki mikserle (biri $\cos$, biri $\sin$ ile) çarpıp $I$ ve $Q$ akışlarını üretir.
Yazılıma ulaşan şey artık taşıyıcıdan arındırılmış, \textbf{karmaşık bant
temeli} (complex baseband) işaretidir: $0$~Hz civarına indirilmiş, hem pozitif
hem negatif frekansları temsil edebilen bir \code{complex<float>} akışı.}

\gsub{Analitik İşaret ve Hilbert Dönüşümü}

Elimizde yalnızca gerçek bir işaret varsa (örneğin bir ses kaydı), onun
karmaşık ``analitik'' karşılığını üretebiliriz. Analitik işaret, negatif
frekansları sıfırlanmış hâlidir: $x_a[n] = x[n] + j\,\hat{x}[n]$, burada
$\hat{x}$ işaretin \textbf{Hilbert dönüşümüdür} ($90^\circ$ faz kaymış hâli).
Bunun en temiz yolu FFT üzerinden gider:

\begin{lstlisting}
#include <complex>
#include <vector>

using cf64 = std::complex<double>;

// Gercek bir isaretin analitik (karmasik) karsiligini uret.
// Negatif frekanslari sifirlayip pozitifleri 2'yle carparak.
std::vector<cf64> analyticSignal(const std::vector<double>& x) {
    const size_t N = x.size();
    std::vector<cf64> X(N);
    for (size_t n = 0; n < N; ++n) X[n] = cf64(x[n], 0.0);

    fft(X);                                   // frekans alanina gec

    // Hilbert filtresi: DC ve Nyquist ayni; pozitif frekanslar x2; negatifler 0
    for (size_t k = 1; k < N/2; ++k)         X[k]     *= 2.0;
    for (size_t k = N/2 + 1; k < N; ++k)     X[k]      = 0.0;

    fft(X, /*invert=*/true);                  // zamana geri don
    return X;                                 // artik analitik (karmasik)
}
\end{lstlisting}

\ipucu{Analitik işaretin büyüklüğü $|x_a[n]|$ işaretin \textbf{anlık zarfını}
(envelope), açısı ise \textbf{anlık fazını} verir. Zarf, AM demodülasyonunun;
fazın türevi ise FM demodülasyonunun temelidir. Yani bir kez karmaşık bant
temeline geçtiğinizde, tüm demodülatörler birkaç satır kalır.}

% =====================================================================
\gsec{NCO, Mikser ve Frekans Kaydırma}
% =====================================================================

Bir işareti frekans ekseninde kaydırmak (mikslemek), onu bir karmaşık üstel ile
çarpmaktır. Bu üsteli üreten yapıya \textbf{NCO} (Numerically Controlled
Oscillator) denir. Bir işareti $-f_0$ kadar kaydırmak (yani $f_0$'daki bir
kanalı DC'ye indirmek) için $e^{-j 2\pi f_0 n / f_s}$ ile çarparız:
\[
y[n] = x[n]\cdot e^{-j 2\pi f_0 n / f_s}
\]

\begin{lstlisting}
#include <complex>
#include <vector>
#include <cmath>

using cf32 = std::complex<float>;

// Sayisal kontrollu osilator (NCO). Her cagrida bir sonraki fazor ornegini verir.
// Faz'i tamsayi bir birikimle degil, kayan noktali ama sarmalanmis tutar.
class Nco {
    double phase_ = 0.0;         // anlik faz (rad)
    double phaseInc_;            // ornek basina faz artisi
public:
    Nco(double freqHz, double fs)
        : phaseInc_(2.0 * M_PI * freqHz / fs) {}

    void setFrequency(double freqHz, double fs) {
        phaseInc_ = 2.0 * M_PI * freqHz / fs;
    }

    cf32 next() {
        cf32 out(std::cos(phase_), std::sin(phase_));
        phase_ += phaseInc_;
        // Fazi [-pi, pi] araliginda tut (sayisal hata birikimini onle)
        if (phase_ >  M_PI) phase_ -= 2.0 * M_PI;
        if (phase_ < -M_PI) phase_ += 2.0 * M_PI;
        return out;
    }
};

// Bir karmasik bant temeli akisini freqHz kadar kaydir (mikser).
void frequencyShift(std::vector<cf32>& iq, double freqHz, double fs) {
    Nco lo(freqHz, fs);          // negatif verirsen asagi kaydirir
    for (cf32& s : iq) s *= lo.next();
}
\end{lstlisting}

\ornek[Kanal Seçimi]{Diyelim ki 2.4~MHz genişlikte bir IQ akışında, merkezden
$+300$~kHz ötede bir istasyon var. Onu tam ortaya (DC'ye) getirmek için akışı
$-300$~kHz kaydırırız; sonra alçak geçiren filtreleyip yalnız o kanalı bırakırız.}

\begin{lstlisting}
void tuneToChannel(std::vector<cf32>& iq, double offsetHz, double fs) {
    frequencyShift(iq, -offsetHz, fs);    // istenen kanali DC'ye getir
    // Sonrasinda: alcak geciren filtre + desimasyon (asagidaki bolum)
}
\end{lstlisting}

\notk[Verimlilik]{Her örnekte \code{sin}/\code{cos} çağırmak pahalıdır. Üretim
kodunda NCO ya bir arama tablosu (LUT) kullanır ya da bir başlangıç fazörünü her
adımda sabit bir \code{cf32} çarpanıyla döndürür: \code{osc *= step;}. İkincisi
çok hızlıdır ama fazörün büyüklüğü zamanla sürüklenir, bu yüzden ara sıra
\code{osc /= std::abs(osc)} ile yeniden normalize edilir.}

% =====================================================================
\gsec{Desimasyon ve Örnekleme Hızı Dönüşümü}
% =====================================================================

Bir SDR akışı genelde 2--10~MHz hızındadır ama biz sonunda 48~kHz'lik ses
istiyoruz. \textbf{Desimasyon} (decimation), örnekleme hızını $M$ kat düşürmektir:
önce alçak geçiren filtre (yeni Nyquist üstünü at, yoksa aliasing!), sonra her
$M$ örnekten birini tut.

\begin{lstlisting}
#include <vector>

// M kat desimasyon: once anti-alias filtre, sonra her M'inci ornegi al.
class Decimator {
    FirFilter aa_;               // anti-aliasing alcak geciren (onceki bolum)
    size_t M_, phase_ = 0;
public:
    Decimator(size_t M, std::vector<float> taps)
        : aa_(std::move(taps)), M_(M) {}

    // Girdi bloklarini isle; cikti girdinin ~1/M kadar orneginidir.
    std::vector<float> process(const std::vector<float>& in) {
        std::vector<float> out;
        out.reserve(in.size() / M_ + 1);
        for (float x : in) {
            float y = aa_.process(x);        // her ornek filtreden gecer
            if (phase_ == 0) out.push_back(y);
            phase_ = (phase_ + 1) % M_;      // ama yalniz 1/M'i saklanir
        }
        return out;
    }
};
\end{lstlisting}

\bilgi{Verimlilik ipucu: Filtre çıktısının çoğunu atacaksak, onları hesaplamak
israftır. \emph{Polifaz desimasyon} yalnızca tutulacak örnekleri hesaplayarak
$M$ kat hızlanma sağlar. Ayrıca çok büyük desimasyon oranları için (örn. $M=100$)
çarpma içermeyen \textbf{CIC filtreleri} kullanılır --- bunlar SDR donanımının
FPGA'sında sıkça bulunur.}

\gsub{Karmaşık (IQ) Desimasyon}

Yukarıdaki \code{FirFilter} gerçek örnekler içindi. IQ için I ve Q kanallarını
ayrı ayrı filtreleriz. Pratikte tek bir karmaşık FIR yazmak daha temizdir:

\begin{lstlisting}
#include <complex>
#include <vector>

using cf32 = std::complex<float>;

// Karmasik girdi, gercek katsayili FIR (kanal filtresi icin tipik).
class ComplexFir {
    std::vector<float> taps_;
    std::vector<cf32> buf_;
    size_t pos_ = 0;
public:
    explicit ComplexFir(std::vector<float> taps)
        : taps_(std::move(taps)), buf_(taps_.size(), cf32(0,0)) {}

    cf32 process(cf32 x) {
        buf_[pos_] = x;
        cf32 acc(0, 0);
        size_t idx = pos_;
        for (size_t k = 0; k < taps_.size(); ++k) {
            acc += taps_[k] * buf_[idx];
            idx = (idx == 0) ? taps_.size() - 1 : idx - 1;
        }
        pos_ = (pos_ + 1) % taps_.size();
        return acc;
    }
};
\end{lstlisting}

% #####################################################################
\gpart{5}{Modülasyon ve Demodülasyon}
% #####################################################################

% =====================================================================
\gsec{Genel Bakış: Modülasyon Neden Var?}
% =====================================================================

\textbf{Modülasyon}, taşımak istediğimiz bilgiyi (mesaj işareti) yüksek frekanslı
bir \emph{taşıyıcıya} (carrier) bindirmektir. Buna neden ihtiyaç var?

\begin{itemize}
\item \textbf{Anten boyutu:} Verimli bir anten dalga boyunun çeyreği kadardır.
1~kHz'lik ses için bu $\sim$75~km'lik bir anten demektir --- imkânsız. 100~MHz
taşıyıcıya bindirirsek anten $\sim$75~cm olur.
\item \textbf{Çoklama:} Farklı istasyonlar farklı taşıyıcı frekanslarını
kullanarak aynı havayı paylaşabilir (frekans bölmeli çoklama).
\item \textbf{Gürültü dayanıklılığı:} FM gibi açı modülasyonları, genlik
gürültüsüne karşı doğal bir bağışıklık sağlar.
\end{itemize}

Bilgiyi taşıyıcının hangi özelliğine gömdüğümüze göre üç temel aile vardır:

\begin{center}
\begin{tabularx}{\textwidth}{@{}>{\bfseries\color{anaMavi}}l >{\raggedright\arraybackslash}X >{\raggedright\arraybackslash}X@{}}
\toprule
\rowcolor{acikMavi}\color{anaMavi}\textbf{Aile} & \color{anaMavi}\textbf{Değişen özellik} & \color{anaMavi}\textbf{Örnekler}\\
\midrule
Genlik & taşıyıcı genliği $A(t)$ & AM, DSB, SSB, ASK, QAM\\
Frekans & anlık frekans $f(t)$ & FM, FSK\\
Faz & anlık faz $\phi(t)$ & PM, BPSK, QPSK, PSK\\
\bottomrule
\end{tabularx}
\end{center}

Karmaşık bant temelinde bir modüle işaret genel olarak şöyle yazılır:
\[
s(t) = A(t)\, e^{\,j\phi(t)}
\]
Genlik modülasyonu $A(t)$'yi, açı modülasyonları (FM/PM) $\phi(t)$'yi bilgiyle
oynar. Bütün demodülatörlerin işi, bu karmaşık işaretten $A(t)$ veya $\phi(t)$'yi
geri çekmektir --- ki bunu Kısım 4'teki \code{std::abs} ve \code{std::arg} ile
yapacağımızı zaten sezdiniz.

% =====================================================================
\gsec{AM: Genlik Modülasyonu}
% =====================================================================

En eski ve en basit modülasyon. Mesaj $m(t)$, taşıyıcının genliğini oynatır:
\[
s(t) = \big(1 + \mu\, m(t)\big)\cos(2\pi f_c t)
\]
Buradaki $\mu$ \textbf{modülasyon indeksidir} (genelde $\le 1$; aşarsa
``aşırı modülasyon'' ve bozulma olur). ``$1+$'' terimi taşıyıcının kendisini de
gönderir --- bu, alıcıda basit bir zarf detektörü kullanmayı mümkün kılar
(bu yüzden eski radyolar ucuzdu).

\gsubsub{AM'in Spektrumu ve Verimliliği}

Tek tonlu $m(t) = \cos(2\pi f_m t)$ için çarpım terimini açalım. Trigonometrik
çarpım-toplam özdeşliği $\cos A \cos B = \tfrac12[\cos(A{-}B) + \cos(A{+}B)]$'yi
kullanırsak:
\[
s(t) = \underbrace{\cos(2\pi f_c t)}_{\text{taşıyıcı}}
     + \underbrace{\tfrac{\mu}{2}\cos\!\big(2\pi (f_c{-}f_m) t\big)}_{\text{alt yan bant}}
     + \underbrace{\tfrac{\mu}{2}\cos\!\big(2\pi (f_c{+}f_m) t\big)}_{\text{üst yan bant}}
\]
Yani AM spektrumu taşıyıcı $f_c$ ile onun $\pm f_m$ yanındaki iki yan banttan
oluşur. Bant genişliği $B = 2 f_m$'dir --- mesajın iki katı. Bilgi \emph{yalnız}
yan bantlardadır; taşıyıcı hiçbir bilgi taşımaz ama tüm gücün çoğunu tüketir.
$\mu = 1$'de bile güç verimliliği sadece $\tfrac{1}{3}$'tür. Bu israfı gidermek
için taşıyıcı bastırılır (\textbf{DSB-SC}) veya bir yan bant da atılır
(\textbf{SSB}, yarı bant genişliği) --- ama bunlar alıcıda senkron demodülasyon
zorunlu kılar.

\gsub{AM Modülatörü}

\begin{lstlisting}
#include <vector>
#include <cmath>

// Klasik cift yan bantli, tasiyicili AM (DSB-LC).
std::vector<float> amModulate(const std::vector<float>& msg,
                              double carrierHz, double fs, double mu = 0.8) {
    std::vector<float> out(msg.size());
    for (size_t n = 0; n < msg.size(); ++n) {
        double carrier = std::cos(2.0 * M_PI * carrierHz * n / fs);
        out[n] = static_cast<float>((1.0 + mu * msg[n]) * carrier);
    }
    return out;
}
\end{lstlisting}

\gsub{Zarf (Envelope) Detektörü ile Demodülasyon}

En basit AM demodülasyonu işaretin zarfını takip etmektir. Klasik analog devre
diyot + RC filtreydi; sayısalda ise ya analitik işaretin büyüklüğünü ($|x_a|$)
alırız ya da doğrudan tam-dalga doğrultup alçak geçiren filtreleriz:

\begin{lstlisting}
#include <vector>
#include <cmath>

// Yontem 1: Analitik isaretin buyuklugu = anlik zarf (en temiz).
std::vector<float> amEnvelopeDemod(const std::vector<double>& rf) {
    auto analytic = analyticSignal(rf);      // Kisim 4'teki fonksiyon
    std::vector<float> msg(analytic.size());
    for (size_t n = 0; n < analytic.size(); ++n)
        msg[n] = static_cast<float>(std::abs(analytic[n]) - 1.0);  // DC'yi cikar
    return msg;
}

// Yontem 2: Dogrultma + alcak geciren (klasik "diyot detektor" dijital hali).
std::vector<float> amRectifierDemod(const std::vector<float>& rf,
                                    double cutoffHz, double fs) {
    Biquad lp = Biquad::lowpass(cutoffHz, fs);
    std::vector<float> msg(rf.size());
    for (size_t n = 0; n < rf.size(); ++n)
        msg[n] = lp.process(std::fabs(rf[n]));   // tam-dalga dogrultma
    // Kalan DC bileseni (tasiyici) ayrica yuksek geciren ile atilabilir.
    return msg;
}
\end{lstlisting}

\gsub{Senkron (Coherent) AM Demodülasyonu}

Zarf detektörü basittir ama gürültüde ve aşırı modülasyonda bozulur. Daha iyisi,
alıcıda taşıyıcının aynısını üretip işaretle çarpmaktır (product detector).
Bant temeline inmiş IQ akışında bu, sadece \textbf{I bileşenini almaktır}:

\begin{lstlisting}
#include <complex>
#include <vector>

using cf32 = std::complex<float>;

// IQ bant temelinde AM: mesaj, tasiyici fazina hizalanmis I bilesenidir.
// (Taban IQ'nun tasiyicisi zaten DC'ye indirilmis kabul edilir.)
std::vector<float> amCoherentDemod(const std::vector<cf32>& iq) {
    std::vector<float> msg(iq.size());
    for (size_t n = 0; n < iq.size(); ++n)
        msg[n] = iq[n].real();               // I = genlik bilgisi
    return msg;
}
\end{lstlisting}

\uyari{Senkron demodülasyon, alıcının taşıyıcı fazını \emph{tam olarak} bilmesini
gerektirir. Gerçekte taşıyıcı fazı bilinmez ve sürüklenir; bu yüzden onu bir
\textbf{PLL} (Faz Kilitlemeli Döngü, Kısım 5.8) ile takip ederiz. PLL kilitlenene
kadar I ve Q karışır ve ses cılız çıkar. Bu, tüm senkron alıcıların ortak
sorunudur.}

% =====================================================================
\gsec{FM: Frekans Modülasyonu}
% =====================================================================

FM, bilgiyi taşıyıcının \textbf{anlık frekansına} gömer. Anlık frekans, mesajla
orantılı olarak nominal değerinin etrafında sapar:
\[
f_i(t) = f_c + k_f\, m(t)
\]
Faz, frekansın integrali olduğundan, modüle işaret:
\[
s(t) = A\cos\!\Big(2\pi f_c t + 2\pi k_f \!\int_0^t m(\tau)\,d\tau\Big)
\]
$k_f$ \textbf{frekans sapmasını} (deviation) belirler. FM radyo yayınında tepe
sapma $\pm 75$~kHz'tir; dar bant FM (telsizler) $\pm 2.5$~kHz kullanır.

\gsub{FM'in Matematiği: Modülasyon İndeksi, Bessel ve Carson}

Tek tonlu bir mesaj $m(t) = \cos(2\pi f_m t)$ ele alalım. Tepe frekans sapması
$\Delta f = k_f$ olsun. Modüle işaret:
\[
s(t) = A\cos\!\Big(2\pi f_c t + \underbrace{\tfrac{\Delta f}{f_m}}_{\beta}\sin(2\pi f_m t)\Big)
\]
Buradaki boyutsuz $\beta = \Delta f / f_m$ niceliğine \textbf{modülasyon indeksi}
denir ve FM'in karakterini belirler. İşareti açmak için karmaşık gösterime
geçelim ve \textbf{Jacobi--Anger açılımını} kullanalım:
\[
e^{\,j\beta\sin\theta} = \sum_{k=-\infty}^{\infty} J_k(\beta)\, e^{\,jk\theta}
\]
Burada $J_k(\beta)$, birinci tür \textbf{Bessel fonksiyonudur}. Bunu $s(t)$'ye
uygularsak:
\[
s(t) = A\sum_{k=-\infty}^{\infty} J_k(\beta)\, \cos\!\big(2\pi (f_c + k f_m)\, t\big)
\]
Bu çarpıcı bir sonuçtur: \emph{tek} bir ton bile FM'de, taşıyıcının etrafında
$f_m$ aralıklarla dizilmiş \emph{sonsuz} sayıda yan bant üretir. Her yan bandın
genliği $J_k(\beta)$ ile belirlenir. Pratikte $|k| > \beta + 1$ için $J_k$ hızla
sıfıra gider, yani bant genişliği sonlu kabul edilebilir. Bu gözlem
\textbf{Carson kuralını} verir --- FM işaretinin gücünün $\%98$'ini içeren
yaklaşık bant genişliği:
\[
B \approx 2(\Delta f + f_m) = 2 f_m(\beta + 1)
\]
FM radyo için $\Delta f = 75$~kHz ve $f_m = 15$~kHz alırsak
$B \approx 2(75 + 15) = 180$~kHz olur --- FM kanallarının neden 200~kHz
aralıkla yerleştirildiğinin sebebi budur.

\notk[Genlik gürültüsüne bağışıklık]{FM'in bilgisi \emph{frekansta} (yani fazın
türevinde) olduğu için, genlik üzerine binen gürültü demodülatörden geçmez ---
alıcıda önce bir sınırlayıcı (limiter) genliği sabitler. Bu, FM'in AM'e kıyasla
neden çok daha temiz ses verdiğinin matematiksel kökenidir: SNR, modülasyon
indeksinin \emph{karesiyle} ($\propto \beta^2$) iyileşir; buna ``FM kazancı''
denir.}

\gsub{FM Modülatörü}

Anahtar nokta: mesajı \emph{integre edip} faza koymak. Ayrık zamanda integral,
birikimli toplamdır:

\begin{lstlisting}
#include <vector>
#include <cmath>

// FM modulator. deviationHz: tepe frekans sapmasi. Mesaj [-1,1] varsayilir.
std::vector<float> fmModulate(const std::vector<float>& msg,
                              double carrierHz, double deviationHz, double fs) {
    std::vector<float> out(msg.size());
    double phase = 0.0;
    const double dt = 1.0 / fs;
    for (size_t n = 0; n < msg.size(); ++n) {
        double fInst = carrierHz + deviationHz * msg[n];   // anlik frekans
        phase += 2.0 * M_PI * fInst * dt;                  // fazi biriktir
        if (phase >  M_PI) phase -= 2.0 * M_PI;
        if (phase < -M_PI) phase += 2.0 * M_PI;
        out[n] = static_cast<float>(std::cos(phase));
    }
    return out;
}

// Karmasik bant temeli FM (SDR ureticileri boyle yapar): tasiyicisiz.
std::vector<cf32> fmModulateBaseband(const std::vector<float>& msg,
                                     double deviationHz, double fs) {
    std::vector<cf32> out(msg.size());
    double phase = 0.0;
    for (size_t n = 0; n < msg.size(); ++n) {
        phase += 2.0 * M_PI * deviationHz * msg[n] / fs;
        out[n] = cf32(std::cos(phase), std::sin(phase));   // e^(j*phase)
    }
    return out;
}
\end{lstlisting}

\gsub{FM Demodülasyonu: IQ Diskriminatör}

İşte belgenin başında sorduğunuz kısım --- ama artık tüm altyapı elimizde
olduğu için nefis derecede basit. FM'de bilgi, fazın \emph{türevindedir} (çünkü
frekans = fazın türevi). Karmaşık bant temelinde anlık faz $\angle z[n]$'dir;
onun türevi ardışık iki örneğin faz farkıdır. Bunu hesaplamanın en zarif yolu
\textbf{polar diskriminatördür}:
\[
m[n] \propto \angle\big(z[n]\cdot z^*[n-1]\big)
\]
Yani ardışık iki IQ örneğinin, birinin eşleniğiyle çarpımının açısı. Bu çarpım
açıları çıkardığı için ($\angle z[n] - \angle z[n-1]$) tam da istediğimiz faz
farkını verir --- üstelik \code{atan2} sarmalanma (wrap) sorunlarını doğal olarak
halleder:

\begin{lstlisting}
#include <complex>
#include <vector>

using cf32 = std::complex<float>;

// Polar diskriminator FM demodulatoru. Girdi: karmasik bant temeli IQ.
// Cikti: demodule mesaj (frekans sapmasiyla orantili).
class FmDemod {
    cf32 prev_{1.0f, 0.0f};      // z[n-1]
    float gain_;
public:
    // gain, cikti seviyesini normalize eder: fs / (2*pi*deviation)
    explicit FmDemod(double fs, double deviationHz)
        : gain_(static_cast<float>(fs / (2.0 * M_PI * deviationHz))) {}

    float process(cf32 z) {
        cf32 prod = z * std::conj(prev_);        // aci farki = arg(prod)
        prev_ = z;
        return std::arg(prod) * gain_;           // atan2 ile faz farki
    }

    std::vector<float> process(const std::vector<cf32>& iq) {
        std::vector<float> out(iq.size());
        for (size_t n = 0; n < iq.size(); ++n) out[n] = process(iq[n]);
        return out;
    }
};
\end{lstlisting}

\ornek[Neden \code{atan2} değil de eşlenik çarpımı?]{Doğrudan
$\angle z[n] - \angle z[n-1]$ hesaplamak da mümkündür ama faz $\pm\pi$
sınırından geçtiğinde $2\pi$'lik sıçramalar oluşur; bunları elle ``unwrap''
etmek gerekir. Eşlenik çarpım yöntemi bu farkı \emph{tek bir} \code{atan2}
çağrısıyla, $[-\pi, \pi]$ aralığında, sıçramasız verir. Bu yüzden GNU Radio'nun
\code{quadrature\_demod} bloğu da tam olarak bunu yapar.}

\notk[Hesaplamadan kaçınmak]{Her örnekte \code{std::arg} (yani \code{atan2})
çağırmak pahalıdır. Yüksek hızlı alıcılarda bunun yerine yaklaşık formül
kullanılır: küçük faz farkları için $\angle(I+jQ) \approx Q/|z|$, hatta çarpım
sonucunun $\frac{I\,\Delta Q - Q\,\Delta I}{I^2+Q^2}$ biçimi. Ama başlangıç için
\code{std::arg} hem doğru hem okunurdur --- önce çalıştırın, sonra hızlandırın.}

\gsub{De-emphasis: FM'in Son Rötuşu}

FM yayıncılığında, verici tarafta yüksek frekanslar bilerek yükseltilir
(pre-emphasis) çünkü gürültü orada daha baskındır. Alıcıda bunu tersine çeviren
bir \textbf{de-emphasis} filtresi (basit bir RC alçak geçiren) uygularız. ABD'de
zaman sabiti 75~µs, Avrupa'da 50~µs'tir:

\begin{lstlisting}
#include <cmath>

// Tek kutuplu de-emphasis filtresi (RC alcak geciren).
class Deemphasis {
    float alpha_, prev_ = 0.0f;
public:
    // tauSec: 75e-6 (ABD) veya 50e-6 (Avrupa/Turkiye)
    Deemphasis(double tauSec, double fs) {
        double dt = 1.0 / fs;
        alpha_ = static_cast<float>(dt / (tauSec + dt));
    }
    float process(float x) {
        prev_ += alpha_ * (x - prev_);           // ussel yumusatma
        return prev_;
    }
};
\end{lstlisting}

% =====================================================================
\gsec{PM: Faz Modülasyonu}
% =====================================================================

Faz modülasyonu, mesajı doğrudan faza koyar (integre etmeden):
\[
s(t) = A\cos\!\big(2\pi f_c t + k_p\, m(t)\big)
\]
FM ile yakın akrabadır: PM'e $m(t)$'nin türevini verirsek FM elde ederiz, ya da
FM'e integralini verirsek PM. Demodülasyon FM'e çok benzer ama faz farkını
\emph{integre etmeyiz} --- doğrudan anlık fazı okuruz:

\begin{lstlisting}
#include <complex>
#include <vector>

using cf32 = std::complex<float>;

// PM modulator (bant temeli): mesaj dogrudan faz.
std::vector<cf32> pmModulateBaseband(const std::vector<float>& msg, float kp) {
    std::vector<cf32> out(msg.size());
    for (size_t n = 0; n < msg.size(); ++n) {
        float phase = kp * msg[n];
        out[n] = cf32(std::cos(phase), std::sin(phase));
    }
    return out;
}

// PM demodulator: anlik fazi oku (unwrap gerekebilir).
std::vector<float> pmDemod(const std::vector<cf32>& iq, float kp) {
    std::vector<float> msg(iq.size());
    float prevPhase = 0.0f;
    for (size_t n = 0; n < iq.size(); ++n) {
        float phase = std::arg(iq[n]);
        // Basit unwrap: 2*pi sicramalarini duzelt
        while (phase - prevPhase >  M_PI) phase -= 2.0f * M_PI;
        while (phase - prevPhase < -M_PI) phase += 2.0f * M_PI;
        prevPhase = phase;
        msg[n] = phase / kp;
    }
    return msg;
}
\end{lstlisting}

% =====================================================================
\gsec{Sayısal Modülasyonlar: ASK, FSK, PSK, QAM}
% =====================================================================

Şimdiye kadar analog mesajları taşıdık. Sayısal haberleşmede ise bit dizileri
taşırız. Bitleri \textbf{sembollere} eşleyip her sembolü karmaşık düzlemde bir
noktaya (\textbf{takım yıldızı / constellation}) koyarız.

\begin{center}
\begin{tabularx}{\textwidth}{@{}>{\bfseries\color{anaMavi}}l >{\raggedright\arraybackslash}X >{\raggedright\arraybackslash}X@{}}
\toprule
\rowcolor{acikMavi}\color{anaMavi}\textbf{Şema} & \color{anaMavi}\textbf{Neyi oynatır} & \color{anaMavi}\textbf{Sembol başına bit}\\
\midrule
ASK / OOK & genlik (var/yok) & 1\\
BPSK & faz (0 / $\pi$) & 1\\
QPSK & faz (4 durum) & 2\\
FSK & frekans (iki ton) & 1\\
16-QAM & genlik + faz (16 nokta) & 4\\
\bottomrule
\end{tabularx}
\end{center}

\gsub{BPSK ve QPSK: Faz Kaydırmalı Anahtarlama}

BPSK, biti taşıyıcının fazına koyar: bit 0 $\to$ $0^\circ$, bit 1 $\to$ $180^\circ$.
QPSK ise her sembolde 2 bit taşır ($4$ faz durumu). Karmaşık bant temelinde bu,
sembol noktasını doğrudan üretmek kadar kolaydır:

\begin{lstlisting}
#include <complex>
#include <vector>
#include <cstdint>

using cf32 = std::complex<float>;

// QPSK: 2 biti bir karmasik sembole esle (Gray kodlu).
// 00 -> ( 1, 1), 01 -> (-1, 1), 11 -> (-1,-1), 10 -> ( 1,-1), 1/sqrt(2) olcekli
cf32 qpskMap(uint8_t twoBits) {
    const float s = 0.70710678f;   // 1/sqrt(2), birim guce normalize
    switch (twoBits & 0x3) {
        case 0b00: return { s,  s};
        case 0b01: return {-s,  s};
        case 0b11: return {-s, -s};
        case 0b10: return { s, -s};
    }
    return {0, 0};
}

// Bit dizisini QPSK sembol dizisine cevir.
std::vector<cf32> qpskModulate(const std::vector<uint8_t>& bits) {
    std::vector<cf32> syms;
    for (size_t i = 0; i + 1 < bits.size(); i += 2) {
        uint8_t sym = (bits[i] << 1) | bits[i+1];
        syms.push_back(qpskMap(sym));
    }
    return syms;
}

// QPSK karar (demod): en yakin takim yildizi noktasina yuvarla.
uint8_t qpskDemap(cf32 r) {
    uint8_t hi = (r.real() < 0) ? 1 : 0;   // I isareti -> ust bit
    uint8_t lo = (r.imag() < 0) ? 1 : 0;   // Q isareti -> alt bit
    // Gray esleme geri cozumu
    uint8_t b0 = hi;
    uint8_t b1 = hi ^ lo;
    return (b0 << 1) | b1;
}
\end{lstlisting}

\gsub{FSK: Frekans Kaydırmalı Anahtarlama}

FSK, her biti farklı bir tonla gönderir. Aslında sayısal bir FM'dir. Modülatörü
FM'inkiyle aynı mantıkta, demodülasyonu ise ya iki band-geçiren filtrenin
enerjisini karşılaştırarak (koherent olmayan) ya da doğrudan FM diskriminatörüyle
yapılır:

\begin{lstlisting}
#include <vector>
#include <cmath>

// Ikili FSK modulator. bit 0 -> f0, bit 1 -> f1. sps: sembol basina ornek.
std::vector<float> fskModulate(const std::vector<uint8_t>& bits,
                               double f0, double f1, double fs, size_t sps) {
    std::vector<float> out;
    double phase = 0.0;
    for (uint8_t b : bits) {
        double f = b ? f1 : f0;
        for (size_t k = 0; k < sps; ++k) {
            phase += 2.0 * M_PI * f / fs;    // faz surekliligi korunur (CPFSK)
            out.push_back(static_cast<float>(std::cos(phase)));
        }
    }
    return out;
}

// Noncoherent FSK demod: iki tonun enerjisini karsilastir (Goertzel benzeri).
std::vector<uint8_t> fskDemod(const std::vector<float>& rf,
                              double f0, double f1, double fs, size_t sps) {
    std::vector<uint8_t> bits;
    for (size_t i = 0; i + sps <= rf.size(); i += sps) {
        double i0=0,q0=0,i1=0,q1=0;
        for (size_t k = 0; k < sps; ++k) {
            double t = (i + k) / fs, x = rf[i + k];
            i0 += x*std::cos(2*M_PI*f0*t);  q0 += x*std::sin(2*M_PI*f0*t);
            i1 += x*std::cos(2*M_PI*f1*t);  q1 += x*std::sin(2*M_PI*f1*t);
        }
        double e0 = i0*i0 + q0*q0, e1 = i1*i1 + q1*q1;
        bits.push_back(e1 > e0 ? 1 : 0);     // hangi ton daha guclu?
    }
    return bits;
}
\end{lstlisting}

\gsub{QAM: Genlik + Faz Birlikte}

QAM (Quadrature Amplitude Modulation), hem genliği hem fazı kullanarak düzlemde
bir ızgara oluşturur. 16-QAM her sembolde 4 bit taşır; WiFi ve DVB gibi yüksek
hızlı sistemlerin temelidir:

\begin{lstlisting}
#include <complex>
#include <vector>
#include <cstdint>

using cf32 = std::complex<float>;

// 16-QAM: 4 biti bir noktaya esle. I ve Q ayri ayri {-3,-1,+1,+3}.
cf32 qam16Map(uint8_t nibble) {
    auto level = [](uint8_t twoBits) -> float {
        switch (twoBits & 0x3) {          // Gray: 00->-3 01->-1 11->+1 10->+3
            case 0b00: return -3.0f; case 0b01: return -1.0f;
            case 0b11: return  1.0f; case 0b10: return  3.0f;
        } return 0;
    };
    float I = level((nibble >> 2) & 0x3);
    float Q = level(nibble & 0x3);
    return cf32(I, Q) / 3.16227766f;      // ortalama guce normalize (~sqrt(10))
}

// En yakin komsu karariyla 16-QAM demod (basit dilim/slicer).
uint8_t qam16Demap(cf32 r) {
    r *= 3.16227766f;                     // olcegi geri al
    auto slice = [](float v) -> uint8_t {
        if (v < -2) return 0b00; if (v < 0) return 0b01;
        if (v < 2)  return 0b11; return 0b10;
    };
    return (slice(r.real()) << 2) | slice(r.imag());
}
\end{lstlisting}

\bilgi{Takım yıldızında noktalar birbirine ne kadar yakınsa, gürültüde karışma
olasılığı o kadar artar. Bu yüzden yüksek dereceli QAM (64, 256, 1024-QAM) daha
fazla bit/sembol taşır ama daha yüksek sinyal/gürültü oranı (SNR) gerektirir.
Modern sistemler kanal kalitesine göre şemayı uyarlar (adaptive modulation):
sinyal iyiyken 256-QAM, kötüyken QPSK'ye düşer.}

% =====================================================================
\gsec{Darbe Şekillendirme (Pulse Shaping)}
% =====================================================================

Sembolleri kutu darbe olarak göndermek, spektrumda geniş yan loblar üretir ve
komşu kanallara taşar. Ayrıca alıcıda \textbf{semboller arası girişim} (ISI)
oluşur. Çözüm, sembolleri yumuşak bir darbeyle şekillendirmektir. Standart seçim
\textbf{kök yükseltilmiş kosinüs} (Root-Raised-Cosine, RRC) filtresidir --- yarısı
vericide, yarısı alıcıda olacak şekilde bölünür:

\begin{lstlisting}
#include <vector>
#include <cmath>

// RRC filtre katsayilari. beta: rolloff (0..1), sps: sembol basina ornek.
std::vector<float> rrcTaps(double beta, size_t sps, size_t numSymbols) {
    const size_t N = sps * numSymbols + 1;   // tek uzunluk, simetrik
    std::vector<float> h(N);
    const double mid = (N - 1) / 2.0;
    for (size_t i = 0; i < N; ++i) {
        double t = (i - mid) / sps;          // sembol cinsinden zaman
        double val;
        if (std::fabs(t) < 1e-8) {
            val = 1.0 - beta + 4.0 * beta / M_PI;
        } else if (std::fabs(std::fabs(4*beta*t) - 1.0) < 1e-8) {
            // Tekillik noktasi (t = 1/(4*beta))
            double a = (1 + 2/M_PI) * std::sin(M_PI/(4*beta));
            double b = (1 - 2/M_PI) * std::cos(M_PI/(4*beta));
            val = beta / std::sqrt(2.0) * (a + b);
        } else {
            double num = std::sin(M_PI*t*(1-beta))
                       + 4*beta*t*std::cos(M_PI*t*(1+beta));
            double den = M_PI*t*(1 - (4*beta*t)*(4*beta*t));
            val = num / den;
        }
        h[i] = static_cast<float>(val);
    }
    // Enerjiyi normalize et
    double e = 0; for (float c : h) e += c*c;
    for (float& c : h) c /= std::sqrt(e);
    return h;
}
\end{lstlisting}

\ipucu{Vericide sembolleri \code{sps} kat upsample edip (araya sıfır koyup) RRC
ile filtrelersiniz. Alıcıda \emph{aynı} RRC ile eşleşmiş filtreleme (matched
filter) yaparsınız. İki RRC'nin kaskadı bir ``yükseltilmiş kosinüs'' verir ki bu
Nyquist ISI kriterini sağlar: sembol anlarında komşu sembollerin katkısı tam
sıfırdır.}

% =====================================================================
\gsec{Taşıyıcı ve Zamanlama Kurtarma: PLL, Costas, Timing}
% =====================================================================

Gerçek bir alıcıda vericinin taşıyıcı frekansı/fazı ve sembol zamanlaması tam
bilinmez. Bunları işaretin kendisinden kestirmek gerekir. Bu senkronizasyon
döngüleri, çalışan bir modemin en zor ve en kritik parçasıdır.

\gsub{Faz Kilitlemeli Döngü (PLL)}

Bir PLL, yerel bir osilatörü gelen işaretin fazına kilitler. Faz dedektörü, döngü
filtresi ve NCO'dan oluşur --- her demodülatörün kalbindeki geri besleme döngüsü:

\begin{lstlisting}
#include <complex>
#include <cmath>

using cf32 = std::complex<float>;

// Karmasik bir tona kilitlenen basit ikinci dereceden PLL.
class Pll {
    double phase_ = 0.0, freq_ = 0.0;
    double alpha_, beta_;                    // dongu kazanclari
public:
    // loopBw: dongu bant genisligi (ornekleme hizina gore normalize, ~0.01)
    explicit Pll(double loopBw) {
        double damping = 0.707;
        double denom = 1 + 2*damping*loopBw + loopBw*loopBw;
        alpha_ = 4*damping*loopBw / denom;   // orantili kazanc
        beta_  = 4*loopBw*loopBw   / denom;  // integral kazanci
    }

    // Girdiyi tasiyicidan arindir ve fazi kilitle. Dondurdugu: duzeltilmis IQ.
    cf32 process(cf32 in) {
        cf32 nco(std::cos(-phase_), std::sin(-phase_));
        cf32 out = in * nco;                 // yerel osilator ile carp
        double err = std::arg(out);          // faz hatasi (dedektor)
        freq_  += beta_  * err;              // frekansi guncelle (integral)
        phase_ += freq_ + alpha_ * err;      // fazi guncelle (orantili)
        if (phase_ >  M_PI) phase_ -= 2*M_PI;
        if (phase_ < -M_PI) phase_ += 2*M_PI;
        return out;
    }
    double frequency() const { return freq_; }
};
\end{lstlisting}

\gsub{Costas Döngüsü: PSK için Taşıyıcı Kurtarma}

PSK işaretlerinde taşıyıcı bastırılmıştır (ayrı bir taşıyıcı tepesi yok), bu
yüzden düz PLL kilitlenemez. \textbf{Costas döngüsü}, faz hatasını modülasyona
duyarsız biçimde üretir. BPSK için hata $I\cdot Q$, QPSK için ise karar-yönlü
bir formdur:

\begin{lstlisting}
#include <complex>
#include <cmath>

using cf32 = std::complex<float>;

// BPSK/QPSK icin Costas dongusu. Tasiyici faz/frekans ofsetini duzeltir.
class CostasLoop {
    double phase_ = 0.0, freq_ = 0.0, alpha_, beta_;
    bool qpsk_;
public:
    CostasLoop(double loopBw, bool qpsk = false) : qpsk_(qpsk) {
        double d = 0.707;
        double denom = 1 + 2*d*loopBw + loopBw*loopBw;
        alpha_ = 4*d*loopBw / denom;
        beta_  = 4*loopBw*loopBw / denom;
    }

    cf32 process(cf32 in) {
        cf32 nco(std::cos(-phase_), std::sin(-phase_));
        cf32 y = in * nco;                   // fazi duzeltilmis ornek

        // Faz hatasi: BPSK -> I*Q, QPSK -> karar yonlu
        double err;
        if (!qpsk_) {
            err = y.real() * y.imag();
        } else {
            float di = (y.real() > 0) ? 1.f : -1.f;   // I karari
            float dq = (y.imag() > 0) ? 1.f : -1.f;   // Q karari
            err = dq * y.real() - di * y.imag();
        }
        freq_  += beta_  * err;
        phase_ += freq_ + alpha_ * err;
        if (phase_ >  M_PI) phase_ -= 2*M_PI;
        if (phase_ < -M_PI) phase_ += 2*M_PI;
        return y;
    }
};
\end{lstlisting}

\gsub{Sembol Zamanlama Kurtarma (Gardner)}

Alıcı, her sembolü tam ortasından örneklemelidir; yoksa ISI ve hatalar artar.
\textbf{Gardner zamanlama dedektörü} örnekleme anını ayarlar ve simge hızının
tam bilinmesini gerektirmez (karar-öncesi çalışır):

\begin{lstlisting}
#include <complex>
#include <vector>

using cf32 = std::complex<float>;

// Gardner TED ile sembol zamanlama kurtarma (sps ornek/sembol, tipik 2+).
// Basitlestirilmis: her sembol icin orta ve tam noktalari kullanir.
std::vector<cf32> gardnerSync(const std::vector<cf32>& x, double sps,
                              double loopGain = 0.01) {
    std::vector<cf32> syms;
    double mu = 0.0;                         // kesirli ornek konumu
    double idx = 0.0;
    cf32 prevSym(0,0), prevMid(0,0);
    while (idx + sps < x.size() - 1) {
        size_t i = static_cast<size_t>(idx);
        cf32 cur = x[i];                     // (pratikte lineer/polifaz interp.)
        cf32 mid = x[static_cast<size_t>(idx - sps/2)];
        // Gardner hatasi: orta ornek, iki sembol gecisini "yakalar"
        double e = (mid.real() * (cur.real() - prevSym.real())
                  + mid.imag() * (cur.imag() - prevSym.imag()));
        mu -= loopGain * e;                  // zamanlamayi ayarla
        syms.push_back(cur);
        prevSym = cur;
        idx += sps + mu;                     // bir sonraki sembole ilerle
        mu *= 0.9;                           // hatayi sonumle
    }
    return syms;
}
\end{lstlisting}

\uyari{Bu senkronizasyon döngüleri pedagojik olarak sadeleştirilmiştir; kesirli
gecikme için gerçek bir interpolatör (Farrow/polifaz), döngü kazançları için
dikkatli ayar ve kilitlenme tespiti gerçek sistemlerde şarttır. Yine de temel
geri besleme yapısı --- \emph{hata üret, filtrele, düzelt} --- her birinde
aynıdır; bu iskeleti kavramak yeter.}

% #####################################################################
\gpart{6}{Uçtan Uca: Tam Bir FM Radyo Alıcısı}
% #####################################################################

% =====================================================================
\gsec{IQ Dosya Formatları ve RTL-SDR}
% =====================================================================

Tüm parçaları bir araya getirip gerçek bir alıcı kuralım. Girdimiz bir SDR'dan
gelen ham IQ akışıdır. En yaygın formatlar:

\begin{center}
\begin{tabularx}{\textwidth}{@{}>{\bfseries\color{anaMavi}}l >{\raggedright\arraybackslash}X@{}}
\toprule
\rowcolor{acikMavi}\color{anaMavi}\textbf{Format} & \color{anaMavi}\textbf{Açıklama}\\
\midrule
\code{cu8} / \code{u8} & RTL-SDR ham çıktısı: I,Q,I,Q\dots\ her biri işaretsiz bayt (0--255).\\
\code{cs8} & işaretli 8-bit (HackRF).\\
\code{cs16} & işaretli 16-bit (Airspy, SDRplay). Yüksek dinamik aralık.\\
\code{cf32} & 32-bit float, $[-1,1]$. GNU Radio'nun iç formatı.\\
\bottomrule
\end{tabularx}
\end{center}

\gsub{RTL-SDR Baytlarını Karmaşık Bant Temeline Çevirme}

\begin{lstlisting}
#include <complex>
#include <vector>
#include <cstdint>
#include <fstream>

using cf32 = std::complex<float>;

// RTL-SDR cu8 baytlarini normalize karmasik ornege cevir.
// 0..255 -> [-1, 1], orta nokta 127.5.
std::vector<cf32> loadRtlSdrIq(const std::string& path) {
    std::ifstream f(path, std::ios::binary);
    std::vector<uint8_t> raw((std::istreambuf_iterator<char>(f)),
                              std::istreambuf_iterator<char>());
    std::vector<cf32> iq(raw.size() / 2);
    for (size_t i = 0; i < iq.size(); ++i) {
        float I = (raw[2*i]     - 127.5f) / 127.5f;
        float Q = (raw[2*i + 1] - 127.5f) / 127.5f;
        iq[i] = cf32(I, Q);
    }
    return iq;
}
\end{lstlisting}

\ipucu{RTL-SDR'dan komut satırıyla ham IQ kaydetmek için:
\code{rtl\_sdr -f 100000000 -s 2400000 -g 40 kayit.iq} --- bu, 100~MHz merkez
frekansta, 2.4~MHz örnekleme hızında, 40~dB kazançla \code{cu8} formatında bir
dosya üretir. Yukarıdaki fonksiyon tam da bu dosyayı okur.}

% =====================================================================
\gsec{Tam WBFM Alıcı Hattı}
% =====================================================================

Şimdi geniş bant FM (WBFM, yani normal FM radyo) için uçtan uca sinyal işleme
hattını kuralım. Adımlar şunlar:

\begin{enumerate}
\item \textbf{Tune:} İstenen istasyonu merkeze (DC) kaydır (NCO mikser).
\item \textbf{Kanal filtresi + desimasyon:} 2.4~MHz'i $\sim$240~kHz'e indir
(FM kanalı $\pm$100~kHz civarı).
\item \textbf{FM demod:} Polar diskriminatörle sesi çek.
\item \textbf{Ses desimasyonu:} 240~kHz'i 48~kHz'e indir.
\item \textbf{De-emphasis:} 50/75~µs RC filtresi.
\item \textbf{WAV'a yaz} veya ses kartına gönder.
\end{enumerate}

\begin{lstlisting}
#include <complex>
#include <vector>
#include <cmath>

using cf32 = std::complex<float>;

// Onceki bolumlerdeki bilesenleri varsayar:
//   Nco, ComplexFir, FirFilter, FmDemod, Deemphasis, designLowpass, writeWav

class WbfmReceiver {
    double sdrRate_;         // ornegin 2.4e6
    double channelRate_;     // ornegin 240e3  (M1 = 10 desimasyon)
    double audioRate_;       // 48e3           (M2 = 5 desimasyon)
    double tuneOffset_;      // istasyonun merkezden ofseti (Hz)

    Nco lo_;
    ComplexFir channelFilter_;
    FmDemod demod_;
    FirFilter audioFilter_;
    Deemphasis deemph_;
    size_t decim1_, decim2_;
public:
    WbfmReceiver(double sdrRate, double tuneOffsetHz)
      : sdrRate_(sdrRate), channelRate_(sdrRate/10.0), audioRate_(48000.0),
        tuneOffset_(tuneOffsetHz),
        lo_(-tuneOffsetHz, sdrRate),                        // asagi kaydir
        channelFilter_(designLowpass(100e3, sdrRate, 65)),  // FM kanal BW
        demod_(channelRate_, 75e3),                         // sapma 75 kHz
        audioFilter_(designLowpass(15e3, channelRate_, 65)),// ses BW 15 kHz
        deemph_(50e-6, audioRate_),                         // Avrupa 50 us
        decim1_(10), decim2_(5) {}

    // Ham IQ blogunu isle -> [-1,1] ses ornekleri (48 kHz mono).
    std::vector<float> process(const std::vector<cf32>& iq) {
        std::vector<float> audio;
        size_t ph1 = 0, ph2 = 0;
        for (cf32 raw : iq) {
            // 1) Tune: istasyonu DC'ye getir
            cf32 shifted = raw * lo_.next();
            // 2) Kanal filtresi (her ornek), desimasyon 1
            cf32 filt = channelFilter_.process(shifted);
            if (ph1++ % decim1_ != 0) continue;
            // 3) FM demodulasyon (artik channelRate_ hizinda)
            float d = demod_.process(filt);
            // 4) Ses filtresi + desimasyon 2
            float a = audioFilter_.process(d);
            if (ph2++ % decim2_ != 0) continue;
            // 5) De-emphasis
            audio.push_back(deemph_.process(a));
        }
        return audio;
    }
};

int main() {
    // 100.0 MHz'te kayit yapildi; 100.3 MHz istasyonunu dinle -> +300 kHz ofset
    auto iq = loadRtlSdrIq("kayit.iq");
    WbfmReceiver rx(2.4e6, +300e3);
    std::vector<float> audio = rx.process(iq);

    // Normalize ve WAV'a yaz
    float peak = 1e-6f;
    for (float s : audio) peak = std::max(peak, std::fabs(s));
    for (float& s : audio) s *= 0.9f / peak;
    writeWav("radyo.wav", audio, 48000, 1);
    return 0;
}
\end{lstlisting}

\ornek[Stereo FM]{FM stereo yayınında sol+sağ (L+R) ses 0--15~kHz'te
gönderilirken, sol$-$sağ (L$-$R) farkı 38~kHz'lik bir alt taşıyıcıda DSB-SC
olarak modüle edilir. 19~kHz'te de bir pilot ton bulunur. Stereo çözmek için:
demod çıktısından 19~kHz pilotu bir PLL ile yakala, frekansını ikiye katla
(38~kHz), L$-$R'yi bu tonla çarparak alçağa indir, sonra $L = (L{+}R)+(L{-}R)$ ve
$R = (L{+}R)-(L{-}R)$ ile iki kanalı ayır. Yukarıdaki hattı mono bıraktık;
stereo, aynı IQ diskriminatör çıktısına eklenen bir DSP katmanıdır.}

% =====================================================================
\gsec{Performans ve Sonraki Adımlar}
% =====================================================================

\gsub{Optimizasyon Notları}

\begin{itemize}
\item \textbf{atan2'den kaçının:} FM demodda her örnekte \code{std::arg}
çağırmak yerine yaklaşık faz-farkı formülleri veya bir CORDIC kullanın.
\item \textbf{Polifaz filtreler:} Desimasyonda atacağınız örnekleri hiç
hesaplamayın --- $M$ kat hızlanma.
\item \textbf{SIMD:} FIR konvolüsyonu, çarpma-toplama ağırlıklıdır; AVX/NEON ile
4--8 kat hızlanır. Modern derleyiciler \code{-O3 -march=native} ile çoğunu
otomatik yapar.
\item \textbf{Blok işleme:} Örnek örnek değil, blok blok çalışın; önbellek ve
dallanma tahmini çok daha verimli olur.
\item \textbf{FFT tabanlı konvolüsyon:} Uzun FIR'lar için (100+ tap),
overlap-add/overlap-save yöntemiyle FFT üzerinden konvolüsyon $O(N\log N)$
verir.
\end{itemize}

\gsub{Üretimde Kullanılan Kütüphaneler}

Bu belge her şeyi sıfırdan yazdı ki mekanizmayı anlayın. Gerçek projelerde ise
olgun kütüphaneleri tercih edin:

\begin{center}
\begin{tabularx}{\textwidth}{@{}>{\bfseries\color{anaMavi}}l >{\raggedright\arraybackslash}X@{}}
\toprule
\rowcolor{acikMavi}\color{anaMavi}\textbf{Araç} & \color{anaMavi}\textbf{Ne için}\\
\midrule
FFTW / KissFFT & Endüstri standardı, çok hızlı FFT (kendi FFT'niz yerine).\\
Liquid-DSP & C ile yazılmış kapsamlı SDR/modem kütüphanesi (filtre, modem, sync).\\
GNU Radio & Blok tabanlı görsel/akış SDR çatısı; prototipleme için mükemmel.\\
VOLK & Vektörize (SIMD) DSP çekirdekleri; GNU Radio'nun altında çalışır.\\
SoapySDR & Donanım-bağımsız SDR sürücü katmanı (RTL, HackRF, Airspy\dots).\\
\bottomrule
\end{tabularx}
\end{center}

\gsub{Nereye Bakmalı}

Buradan devam etmek için doğal yollar: \textbf{sayısal alıcı zinciri} (AGC,
eşleşmiş filtre, senkron, çözücü), \textbf{kanal kodlaması} (Hamming, konvolüsyon,
Viterbi, LDPC, Turbo), \textbf{OFDM} (WiFi/LTE/5G'nin çok taşıyıcılı temeli),
\textbf{spektrum kestirimi} (Welch, MUSIC), ve \textbf{uyarlamalı filtreler}
(LMS, RLS ile ekolayzasyon ve gürültü giderme).

\bilgi{Öğrenmenin en hızlı yolu \emph{dinlemek} ve \emph{görmektir}. Her
demodülatör aşamasının çıktısını WAV'a yazıp dinleyin, FFT'sini alıp spektrumuna
bakın. Bir FM alıcısını çalıştırıp gerçek bir istasyonu duyduğunuz an, bu
belgedeki tüm parçalar --- örnekleme, IQ, mikser, filtre, diskriminatör --- tek
bir bütün olarak yerine oturur.}

% #####################################################################
\gpart{7}{İleri Konular}
% #####################################################################

% =====================================================================
\gsec{Stereo FM Çözücü (Tam Uygulama)}
% =====================================================================

Kısım 6'da mono FM alıcısını kurmuş, stereo'yu ise kavramsal bırakmıştık. Şimdi
onu tam koda dökelim. FM stereo yayınında demodülatör çıktısı (\textbf{MPX} ya da
kompozit işaret) şu katmanlardan oluşur:

\begin{center}
\begin{tabularx}{\textwidth}{@{}>{\bfseries\color{anaMavi}}l >{\raggedright\arraybackslash}X@{}}
\toprule
\rowcolor{acikMavi}\color{anaMavi}\textbf{Bileşen} & \color{anaMavi}\textbf{Frekans bandı}\\
\midrule
L+R (mono uyumlu) & 0--15 kHz (bant temeli)\\
Pilot ton & 19 kHz (saf sinüs, faz referansı)\\
L$-$R (DSB-SC) & 23--53 kHz, 38 kHz alt taşıyıcı etrafında\\
RDS (opsiyonel) & 57 kHz (pilotun 3. harmoniği)\\
\bottomrule
\end{tabularx}
\end{center}

Sol ve sağ kanalı ayırmak için: (1) 19~kHz pilotu bir PLL ile yakala, (2)
frekansını ikiye katlayıp 38~kHz üret, (3) L$-$R'yi bu tonla çarpıp bant temeline
indir, (4) matris işlemiyle $L = (L{+}R) + (L{-}R)$ ve $R = (L{+}R) - (L{-}R)$
hesapla.

\begin{lstlisting}
#include <complex>
#include <vector>
#include <cmath>

// Kompozit (MPX) FM cikisindan stereo L/R cozer.
// Girdi: FM demod cikisi (Kisim 6'daki FmDemod), ~240 kHz hizinda.
// Cikti: interleaved stereo (L,R,L,R...) 48 kHz'e desime edilmis.
class StereoDecoder {
    double mpxRate_;                 // MPX ornekleme hizi (orn. 240 kHz)
    Pll pilotPll_;                   // 19 kHz pilota kilitlenir (Kisim 5)
    Biquad pilotBpf_;                // 19 kHz band geciren
    FirFilter lprLpf_, lmrLpf_;      // her iki kanal icin 15 kHz alcak geciren
    Deemphasis deL_, deR_;
    size_t decim_, phase_ = 0;
public:
    StereoDecoder(double mpxRate, double audioRate)
      : mpxRate_(mpxRate),
        pilotPll_(0.001),                                  // dar dongu BW
        pilotBpf_(Biquad::bandpass(19000.0, mpxRate, 20)),
        lprLpf_(designLowpass(15000.0, mpxRate, 65)),
        lmrLpf_(designLowpass(15000.0, mpxRate, 65)),
        deL_(50e-6, audioRate), deR_(50e-6, audioRate),
        decim_(static_cast<size_t>(mpxRate / audioRate)) {}

    std::vector<float> process(const std::vector<float>& mpx) {
        std::vector<float> stereo;
        for (float x : mpx) {
            // 1) 19 kHz pilotu izole et ve PLL'e ver
            float pilot = pilotBpf_.process(x);
            // PLL'i gercek pilota kilitle; kilitli fazi al
            cf32 locked = pilotPll_.process(cf32(pilot, 0.0f));
            double ph = std::arg(pilotPll_.phaseState());  // 19 kHz fazi

            // 2) 38 kHz alt tasiyici = pilot fazinin 2 kati
            float sub38 = std::cos(2.0 * ph);

            // 3) Kanallari ayir
            float lpr = lprLpf_.process(x);           // L+R (bant temeli)
            float lmr = lmrLpf_.process(x * sub38 * 2.0f); // L-R (asagi mixle)

            // 4) Desimasyon
            if (phase_++ % decim_ != 0) continue;

            // 5) Matris + de-emphasis
            float L = deL_.process(lpr + lmr);
            float R = deR_.process(lpr - lmr);
            stereo.push_back(L);
            stereo.push_back(R);
        }
        return stereo;
    }
};
\end{lstlisting}

\notk[PLL'e küçük bir ekleme]{Yukarıdaki kod, Kısım 5'teki \code{Pll} sınıfına
anlık fazı dışarı veren bir yardımcı ekler:
\code{cf32 phaseState() const \{ return \{std::cos(phase\_), std::sin(phase\_)\}; \}}
Böylece pilotun fazını okuyup ikiye katlayarak 38~kHz üretiriz. Pilot yoksa
(zayıf sinyal) PLL kilitlenemez; bu durumda yalnız L+R kullanılıp otomatik olarak
mono'ya düşülür --- gerçek alıcıların yaptığı budur.}

\ipucu{Stereo çözmenin başarısı tamamen pilot PLL'in kilidine bağlıdır. 19~kHz'e
$\pm$2~Hz'ten fazla sapma stereo ayrımını (separation) bozar. Bu yüzden döngü
bant genişliğini dar (0.001 civarı) tutarız: yavaş kilitlenir ama gürültüye
karşı çok kararlıdır.}

% =====================================================================
\gsec{Otomatik Kazanç Kontrolü (AGC)}
% =====================================================================

Gerçek sinyallerin gücü sürekli değişir (sönümlenme/fading, mesafe, anten yönü).
Demodülatörlerin çoğu \emph{sabit genlik} bekler --- özellikle QAM gibi genlik
taşıyan şemalar. \textbf{AGC}, işareti hedef bir güç seviyesine sürekli
ölçekleyerek bunu sağlar:

\begin{lstlisting}
#include <complex>
#include <cmath>

using cf32 = std::complex<float>;

// Karmasik AGC: cikis gucunu hedef seviyeye kilitler.
class Agc {
    float gain_ = 1.0f;
    float target_, rate_;
public:
    // target: istenen cikis genligi (orn. 1.0), rate: uyum hizi (0..1, ~1e-3)
    Agc(float target = 1.0f, float rate = 1e-3f)
        : target_(target), rate_(rate) {}

    cf32 process(cf32 x) {
        cf32 out = x * gain_;
        float mag = std::abs(out);
        // Hata: hedef ile olculen arasindaki fark -> kazanci ayarla
        gain_ += rate_ * (target_ - mag) * gain_;
        if (gain_ < 1e-6f) gain_ = 1e-6f;    // patlamayi onle
        return out;
    }
    float gain() const { return gain_; }
};
\end{lstlisting}

\uyari{AGC hızı (\code{rate}) kritik bir dengedir. Çok hızlıysa, işaretin kendi
genlik modülasyonunu (AM/QAM bilgisi!) ``düzleştirip'' bilgiyi siler. Çok
yavaşsa, ani sönümlenmelere yetişemez. Kural: AGC, sembol hızından çok daha yavaş
ama fading'den çok daha hızlı olmalıdır.}

% =====================================================================
\gsec{Uçtan Uca Sayısal Modem: QPSK Vericiden Alıcıya}
% =====================================================================

Şimdi Kısım 5'teki tüm parçaları --- QPSK eşleme, RRC şekillendirme, kanal, AGC,
Costas ve zamanlama --- tek bir çalışan modemde birleştirelim. Bu, gerçek bir
sayısal haberleşme zincirinin iskeletidir.

\gsub{Verici (Modülatör)}

\begin{lstlisting}
#include <complex>
#include <vector>
#include <cstdint>

using cf32 = std::complex<float>;

// Bitler -> QPSK sembol -> upsample -> RRC sekillendirme -> baseband IQ
std::vector<cf32> qpskTransmit(const std::vector<uint8_t>& bits,
                               size_t sps, double beta) {
    // 1) Bitleri QPSK sembollere esle (Kisim 5)
    std::vector<cf32> syms = qpskModulate(bits);

    // 2) sps kat upsample: her sembol arasina sps-1 sifir koy
    std::vector<cf32> up(syms.size() * sps, cf32(0, 0));
    for (size_t i = 0; i < syms.size(); ++i) up[i * sps] = syms[i];

    // 3) RRC filtre ile darbe sekillendir (Kisim 5'teki rrcTaps)
    ComplexFir shaper(rrcTaps(beta, sps, 6));
    std::vector<cf32> tx(up.size());
    for (size_t i = 0; i < up.size(); ++i) tx[i] = shaper.process(up[i]);
    return tx;
}
\end{lstlisting}

\gsub{Kanal Modeli}

Alıcıyı test etmek için gerçekçi bozulmalar ekleyen bir kanal yazalım: frekans
ofseti, faz kayması ve Gauss gürültüsü. (\code{std::mt19937} ile deterministik.)

\begin{lstlisting}
#include <random>

// Kanal: frekans ofseti + faz + AWGN gurultusu ekler.
std::vector<cf32> channel(const std::vector<cf32>& tx, double freqOffset,
                          double phase, double noiseStd, uint32_t seed = 42) {
    std::mt19937 rng(seed);
    std::normal_distribution<float> noise(0.0f, noiseStd);
    std::vector<cf32> rx(tx.size());
    double ph = phase;
    for (size_t i = 0; i < tx.size(); ++i) {
        cf32 rot(std::cos(ph), std::sin(ph));         // tasiyici ofseti
        rx[i] = tx[i] * rot + cf32(noise(rng), noise(rng));
        ph += 2.0 * M_PI * freqOffset;
    }
    return rx;
}
\end{lstlisting}

\gsub{Alıcı (Demodülatör)}

Alıcı zinciri: AGC $\to$ eşleşmiş RRC filtre $\to$ zamanlama kurtarma $\to$
Costas ile taşıyıcı kurtarma $\to$ karar (slicer) $\to$ bitler.

\begin{lstlisting}
#include <complex>
#include <vector>
#include <cstdint>

using cf32 = std::complex<float>;

std::vector<uint8_t> qpskReceive(std::vector<cf32> rx, size_t sps, double beta) {
    // 1) AGC ile genligi normalize et
    Agc agc(1.0f, 1e-3f);
    for (cf32& s : rx) s = agc.process(s);

    // 2) Eslesmis filtre (verici ile ayni RRC) -> ISI'yi minimize eder
    ComplexFir matched(rrcTaps(beta, sps, 6));
    std::vector<cf32> filt(rx.size());
    for (size_t i = 0; i < rx.size(); ++i) filt[i] = matched.process(rx[i]);

    // 3) Sembol zamanlama kurtarma (Kisim 5'teki Gardner) -> sembol hizina in
    std::vector<cf32> syms = gardnerSync(filt, static_cast<double>(sps));

    // 4) Costas dongusu ile tasiyici faz/frekans kurtar (QPSK modu)
    CostasLoop costas(0.01, /*qpsk=*/true);
    std::vector<uint8_t> bits;
    for (cf32 s : syms) {
        cf32 corrected = costas.process(s);
        // 5) Karar: en yakin QPSK noktasi -> 2 bit (Kisim 5'teki qpskDemap)
        uint8_t sym = qpskDemap(corrected);
        bits.push_back((sym >> 1) & 1);
        bits.push_back(sym & 1);
    }
    return bits;
}

// Uctan uca test: gonder -> kanaldan gecir -> al -> bit hata oranini (BER) olc.
#include <cstdio>
int main() {
    std::vector<uint8_t> tx_bits;
    std::mt19937 rng(1);
    for (int i = 0; i < 2000; ++i) tx_bits.push_back(rng() & 1);

    auto tx  = qpskTransmit(tx_bits, /*sps=*/8, /*beta=*/0.35);
    auto rx  = channel(tx, /*freqOff=*/1e-4, /*phase=*/0.7, /*noise=*/0.1);
    auto out = qpskReceive(rx, 8, 0.35);

    // BER: senkron gecikmesi nedeniyle hizalanmis kismi karsilastir
    size_t errors = 0, compared = 0;
    size_t offset = 40;   // filtre/sync gecikmesi (deneysel)
    for (size_t i = 0; i + offset < out.size() && i < tx_bits.size(); ++i) {
        errors += (out[i + offset] != tx_bits[i]);
        ++compared;
    }
    printf("BER = %.4f (%zu/%zu)\n", (double)errors/compared, errors, compared);
    return 0;
}
\end{lstlisting}

\bilgi{Gerçek bir modemde ek olarak: \textbf{çerçeve senkronizasyonu} (bilinen
bir önsöz/preamble dizisiyle bit hizalaması), \textbf{faz belirsizliği çözümü}
(QPSK $90^\circ$'lik dört olası dönüşten hangisinde kilitlendiğini bilmez ---
diferansiyel kodlama veya bilinen semboller çözer) ve \textbf{kanal
denkleştirme} (equalization) bulunur. Yukarıdaki iskelet bu eklemeler için
sağlam bir temeldir.}

% =====================================================================
\gsec{OFDM: Çok Taşıyıcılı Modülasyon}
% =====================================================================

WiFi, LTE, 5G, DVB-T ve DAB'ın temelinde \textbf{OFDM} (Orthogonal
Frequency-Division Multiplexing) yatar. Fikir zarif: yüksek hızlı bir bit
akışını, her biri kendi dar taşıyıcısında düşük hızla giden \emph{binlerce}
paralel alt kanala böl. Bu alt taşıyıcılar dik seçilir, öyle ki bir \textbf{IFFT}
ile hepsi tek adımda üretilir --- OFDM'in tüm sihri budur.

\begin{lstlisting}
#include <complex>
#include <vector>

using cf64 = std::complex<double>;

// OFDM verici: her alt tasiyiciya bir QAM sembolu koy, IFFT ile zamana don,
// arkasina dongusel onek (cyclic prefix) ekle.
std::vector<cf64> ofdmModulate(const std::vector<cf64>& qamSymbols,
                               size_t nfft, size_t cpLen) {
    std::vector<cf64> out;
    for (size_t base = 0; base < qamSymbols.size(); base += nfft) {
        // 1) Bir FFT-blogu doldur (alt tasiyicilar = frekans binleri)
        std::vector<cf64> freq(nfft, cf64(0, 0));
        for (size_t k = 0; k < nfft && base + k < qamSymbols.size(); ++k)
            freq[k] = qamSymbols[base + k];

        // 2) IFFT -> zaman alani OFDM sembolu (Kisim 2'deki fft, invert=true)
        fft(freq, /*invert=*/true);

        // 3) Dongusel onek: son cpLen ornegi basa kopyala
        //    -> cok yollu yayilima (multipath) karsi bagisiklik saglar
        for (size_t i = nfft - cpLen; i < nfft; ++i) out.push_back(freq[i]);
        for (size_t i = 0; i < nfft; ++i)            out.push_back(freq[i]);
    }
    return out;
}

// OFDM alici: dongusel oneki at, FFT ile alt tasiyicilari geri cek.
std::vector<cf64> ofdmDemodulate(const std::vector<cf64>& rx,
                                 size_t nfft, size_t cpLen) {
    std::vector<cf64> symbols;
    const size_t symLen = nfft + cpLen;
    for (size_t base = 0; base + symLen <= rx.size(); base += symLen) {
        // 1) Dongusel oneki atla
        std::vector<cf64> block(rx.begin() + base + cpLen,
                                rx.begin() + base + cpLen + nfft);
        // 2) FFT -> her bin bir alt tasiyicinin QAM sembolu
        fft(block, /*invert=*/false);
        for (cf64 s : block) symbols.push_back(s);
    }
    return symbols;
}
\end{lstlisting}

\notk[Döngüsel önek neden işe yarar?]{Çok yollu bir kanalda işaretin gecikmeli
kopyaları üst üste biner (semboller arası girişim). Döngüsel önek, her OFDM
sembolünün kuyruğunu başına kopyalayarak lineer konvolüsyonu \emph{dairesel}
konvolüsyona çevirir. Dairesel konvolüsyon frekans alanında basit bir
\emph{çarpmaya} dönüştüğü için (konvolüsyon teoremi, Kısım 8), her alt taşıyıcı
kanalı tek bir karmaşık katsayıyla düzeltilebilir --- OFDM'in çok yollu kanallara
karşı efsanevi dayanıklılığı buradan gelir.}

% =====================================================================
\gsec{Kanal Kodlaması: Konvolüsyonel Kod ve Viterbi}
% =====================================================================

Gürültü kaçınılmaz olarak bit hatası üretir. \textbf{İleri hata düzeltme} (FEC),
gönderilen veriye fazlalık ekleyerek alıcının hataları \emph{düzeltmesini}
sağlar. En klasik yöntem, \textbf{konvolüsyonel kod} + \textbf{Viterbi
çözücüdür} (uydu, WiFi, GSM'de kullanılır).

\gsub{Konvolüsyonel Kodlayıcı}

Kısıt uzunluğu $K=3$, oran $1/2$ bir kodlayıcı: her giriş biti için, kaydırma
yazmacındaki son 3 bitten iki çıkış biti üretilir (üreteç çokterlileri
$G_1 = 111_2 = 7$, $G_2 = 101_2 = 5$).

\begin{lstlisting}
#include <vector>
#include <cstdint>

// Oran 1/2, K=3 konvolusyonel kodlayici. Her giris biti -> 2 cikis biti.
std::vector<uint8_t> convEncode(const std::vector<uint8_t>& bits) {
    uint8_t state = 0;                       // 2 bitlik kaydirma yazmaci
    std::vector<uint8_t> out;
    for (uint8_t b : bits) {
        uint8_t reg = ((b & 1) << 2) | state; // 3 bit: [yeni, s1, s0]
        // Parite: G1=7 (111), G2=5 (101) -> XOR ile bit sayimi
        uint8_t o1 = __builtin_parity(reg & 0b111);
        uint8_t o2 = __builtin_parity(reg & 0b101);
        out.push_back(o1);
        out.push_back(o2);
        state = reg & 0b11;                  // yeni durum (son 2 bit)
    }
    return out;
}
\end{lstlisting}

\gsub{Viterbi Çözücü}

Viterbi, tüm olası durum dizileri (trellis) arasından, alınan bitlere en yakın
yolu \emph{en olası gönderilen dizi} olarak bulur --- dinamik programlamayla
verimli biçimde:

\begin{lstlisting}
#include <vector>
#include <cstdint>
#include <limits>
#include <algorithm>

// Oran 1/2, K=3 icin sert-karar (hard-decision) Viterbi cozucu.
std::vector<uint8_t> viterbiDecode(const std::vector<uint8_t>& recv) {
    const int NS = 4;                        // durum sayisi (2^(K-1))
    const int INF = std::numeric_limits<int>::max() / 2;
    size_t steps = recv.size() / 2;

    // Her adimda: her durumun en dusuk maliyeti ve gelinen yol
    std::vector<std::vector<int>> cost(steps + 1, std::vector<int>(NS, INF));
    std::vector<std::vector<int>> prev(steps + 1, std::vector<int>(NS, -1));
    std::vector<std::vector<uint8_t>> inBit(steps + 1, std::vector<uint8_t>(NS));
    cost[0][0] = 0;                          // baslangic durumu 0

    // Ileri gecis: trellis'i doldur (Hamming mesafesi = dal metrigi)
    for (size_t t = 0; t < steps; ++t) {
        uint8_t r1 = recv[2*t], r2 = recv[2*t + 1];
        for (int s = 0; s < NS; ++s) {
            if (cost[t][s] == INF) continue;
            for (uint8_t b = 0; b < 2; ++b) {          // olasi giris biti
                uint8_t reg = (b << 2) | s;
                uint8_t e1 = __builtin_parity(reg & 0b111);
                uint8_t e2 = __builtin_parity(reg & 0b101);
                int metric = (e1 != r1) + (e2 != r2);  // Hamming mesafesi
                int ns = reg & 0b11;
                int nc = cost[t][s] + metric;
                if (nc < cost[t+1][ns]) {
                    cost[t+1][ns] = nc;
                    prev[t+1][ns] = s;
                    inBit[t+1][ns] = b;
                }
            }
        }
    }

    // En dusuk maliyetli son durumu bul, geriye iz surerek bitleri kurtar
    int best = 0;
    for (int s = 1; s < NS; ++s) if (cost[steps][s] < cost[steps][best]) best = s;

    std::vector<uint8_t> decoded(steps);
    int s = best;
    for (size_t t = steps; t > 0; --t) {
        decoded[t-1] = inBit[t][s];
        s = prev[t][s];
    }
    return decoded;
}
\end{lstlisting}

\ornek[Hata düzeltmeyi görmek]{Kodla, birkaç biti bilerek boz, çöz --- Viterbi
onları düzeltir:}

\begin{lstlisting}
#include <cstdio>
int main() {
    std::vector<uint8_t> msg = {1,0,1,1,0,0,1,0,1,1};
    auto enc = convEncode(msg);
    enc[3] ^= 1;  enc[10] ^= 1;              // 2 kanal hatasi enjekte et
    auto dec = viterbiDecode(enc);

    int errs = 0;
    for (size_t i = 0; i < msg.size(); ++i) errs += (dec[i] != msg[i]);
    printf("Kanal hatasi vardi ama cozum sonrasi hata = %d\n", errs); // 0
    return 0;
}
\end{lstlisting}

\ipucu{Yukarıdaki ``sert karar'' Viterbi, bitleri 0/1'e yuvarlanmış alır.
\textbf{Yumuşak karar} (soft-decision) Viterbi ise demodülatörün ham güven
değerlerini (örneğin Costas çıkışındaki mesafeleri) kullanır ve $\sim$2~dB daha
iyi başarım verir. Modern sistemler bunu bir adım öteye taşıyıp \textbf{Turbo} ve
\textbf{LDPC} kodları kullanır --- Shannon sınırına çok yaklaşırlar.}

% #####################################################################
\gpart{8}{Projeler}
% #####################################################################

Bu kısım, önceki bölümlerdeki yapı taşlarını birleştiren dört \emph{tam,
derlenip çalışan} projeyi adım adım kurar. Her proje kendi başına bir dosyadır;
bu belgedeki yardımcıları (WAV G/Ç, FFT, filtreler) yeniden kullanır. En sonda
hepsini derleyen bir CMake dosyası var.

% =====================================================================
\gsec{Proje 1 --- DTMF (Telefon Tuşu) Çözücü}
% =====================================================================

Telefon tuşlarının çıkardığı ``tuu-tuu'' sesleri \textbf{DTMF}'dir (Dual-Tone
Multi-Frequency): her tuş, biri düşük biri yüksek gruptan iki sinüsün toplamıdır.
Örneğin ``5'' tuşu 770~Hz + 1336~Hz'tir. Çözmek için, bu 8 frekansın her birinin
enerjisini ölçüp en güçlü düşük+yüksek çifti buluruz. Tek bir frekansın enerjisini
tüm FFT'yi hesaplamadan bulan \textbf{Goertzel algoritması} tam bu iş içindir.

\begin{lstlisting}
#include <vector>
#include <cmath>
#include <cstdio>
#include <array>

// Goertzel: tek bir hedef frekansin enerjisini O(N)'de, FFT'siz olcer.
double goertzel(const std::vector<float>& x, double targetHz, double fs) {
    const size_t N = x.size();
    double k = 0.5 + (N * targetHz) / fs;
    double w = 2.0 * M_PI * static_cast<int>(k) / N;
    double coeff = 2.0 * std::cos(w);
    double s0 = 0, s1 = 0, s2 = 0;
    for (size_t n = 0; n < N; ++n) {
        s0 = x[n] + coeff * s1 - s2;
        s2 = s1;  s1 = s0;
    }
    return s1*s1 + s2*s2 - coeff*s1*s2;    // enerji (buyukluk karesi)
}

// Bir ses blogundan DTMF tusunu coz. Tanimsizsa 0 doner.
char decodeDtmf(const std::vector<float>& block, double fs) {
    static const double lowF[4]  = {697, 770, 852, 941};
    static const double highF[4] = {1209, 1336, 1477, 1633};
    static const char keys[4][4] = {
        {'1','2','3','A'}, {'4','5','6','B'},
        {'7','8','9','C'}, {'*','0','#','D'}
    };
    // En guclu dusuk ve yuksek grup frekansini bul
    int lo = 0, hi = 0; double bestLo = 0, bestHi = 0;
    for (int i = 0; i < 4; ++i) {
        double eL = goertzel(block, lowF[i],  fs);
        double eH = goertzel(block, highF[i], fs);
        if (eL > bestLo) { bestLo = eL; lo = i; }
        if (eH > bestHi) { bestHi = eH; hi = i; }
    }
    // Esik: iki ton da yeterince guclu mu? (basit sessizlik reddi)
    double avg = 0; for (float v : block) avg += v*v;
    avg /= block.size();
    if (bestLo < avg * block.size() * 0.5 ||
        bestHi < avg * block.size() * 0.5) return 0;
    return keys[lo][hi];
}

int main() {
    const double fs = 8000.0;   // telefon standardi
    // Test: "5" tusunu sentezle (770 + 1336 Hz), sonra coz
    std::vector<float> tone;
    for (int n = 0; n < 400; ++n) {       // 50 ms blok
        double t = n / fs;
        tone.push_back(0.5f*std::sin(2*M_PI*770*t)
                     + 0.5f*std::sin(2*M_PI*1336*t));
    }
    printf("Cozulen tus: %c\n", decodeDtmf(tone, fs));   // '5'
    return 0;
}
\end{lstlisting}

\ipucu{Gerçek bir telefon kaydını çözmek için sesi 40--50~ms'lik bloklara bölüp
her bloğu \code{decodeDtmf}'e verin; ardışık aynı tuşları teke indirin (debounce).
Bir \code{.wav} kaydı okumak için Kısım 1'deki WAV okuyucuyu ekleyin --- gerçek
bir çağrı merkezinin IVR sistemi tam olarak böyle çalışır.}

% =====================================================================
\gsec{Proje 2 --- Ses Üzerinden Veri İletimi (FSK Modem)}
% =====================================================================

Bu proje, iki bilgisayarın hoparlör ve mikrofonla veri alışverişi yapmasını
sağlayan minik bir modem kurar --- ``ggwave'' ve eski telefon modemlerinin
prensibi. Metni bitlere çevirip FSK (Kısım 5) ile ses tonlarına, sonra tekrar
metne dönüştürürüz. Senkronizasyon için başa bilinen bir önsöz (preamble) tonu
koyarız.

\begin{lstlisting}
#include <vector>
#include <string>
#include <cstdint>
#include <cmath>

// --- VERICI: metin -> ses ornekleri ---
std::vector<float> encodeText(const std::string& text, double fs) {
    const double f0 = 1200.0, f1 = 2200.0;   // Bell 202 benzeri tonlar
    const size_t sps = static_cast<size_t>(fs / 300.0);  // 300 baud
    // Metni bit dizisine ac (her karakter 8 bit, MSB once)
    std::vector<uint8_t> bits;
    // Onsoz: senkron icin 16 adet 1-0 gecisi
    for (int i = 0; i < 16; ++i) { bits.push_back(1); bits.push_back(0); }
    for (char c : text)
        for (int b = 7; b >= 0; --b) bits.push_back((c >> b) & 1);

    // FSK modulasyon (surekli faz)
    std::vector<float> out;
    double phase = 0.0;
    for (uint8_t bit : bits) {
        double f = bit ? f1 : f0;
        for (size_t k = 0; k < sps; ++k) {
            phase += 2.0 * M_PI * f / fs;
            out.push_back(0.6f * static_cast<float>(std::sin(phase)));
        }
    }
    return out;
}

// --- ALICI: ses ornekleri -> metin ---
std::string decodeText(const std::vector<float>& audio, double fs) {
    const double f0 = 1200.0, f1 = 2200.0;
    const size_t sps = static_cast<size_t>(fs / 300.0);

    // Her sembol icin iki tonun enerjisini Goertzel ile karsilastir
    std::vector<uint8_t> bits;
    for (size_t i = 0; i + sps <= audio.size(); i += sps) {
        std::vector<float> block(audio.begin() + i, audio.begin() + i + sps);
        double e0 = goertzel(block, f0, fs);   // Proje 1'deki goertzel
        double e1 = goertzel(block, f1, fs);
        bits.push_back(e1 > e0 ? 1 : 0);
    }

    // Onsozden sonraki veriyi bul: 32 bitlik gecis desenini atla
    size_t start = 32;
    // Bitleri baytlara topla
    std::string text;
    for (size_t i = start; i + 8 <= bits.size(); i += 8) {
        uint8_t c = 0;
        for (int b = 0; b < 8; ++b) c = (c << 1) | bits[i + b];
        if (c == 0) break;
        text.push_back(static_cast<char>(c));
    }
    return text;
}

#include <cstdio>
int main() {
    const double fs = 44100.0;
    std::string msg = "MERHABA DSP";
    auto audio = encodeText(msg, fs);
    writeWav("modem.wav", audio, 44100, 1);      // Kisim 1'deki writeWav
    std::string got = decodeText(audio, fs);
    printf("Gonderilen: %s\nAlinan    : %s\n", msg.c_str(), got.c_str());
    return 0;
}
\end{lstlisting}

\bilgi{\code{modem.wav}'ı bir bilgisayarda çalıp diğerinde mikrofonla kaydederek
gerçek havadan iletim yapabilirsiniz. Gerçek ortamda çalışması için ekstra:
bit hatalarına karşı Kısım 7'deki Viterbi kodu, ve önsözden sonra tam bit
hizası için bir kare senkronizasyon damgası (örn. \code{0x7E} bayrağı) ekleyin.}

% =====================================================================
\gsec{Proje 3 --- Terminal Spektrum Analizörü}
% =====================================================================

Bir WAV dosyasını okuyup, kayan pencerelerle FFT alarak terminalde ASCII bir
spektrum (mini ``şelale'') çizen bir araç. Bir sinyalin frekans içeriğini
gözle görmenin en hızlı yolu --- hata ayıklarken paha biçilmezdir.

\begin{lstlisting}
#include <vector>
#include <complex>
#include <cstdio>
#include <cmath>
#include <algorithm>

using cf64 = std::complex<double>;

// [-1,1] mono ornekleri al, ardisik FFT pencerelerini ASCII bar olarak ciz.
void spectrumAnalyzer(const std::vector<float>& x, double fs) {
    const size_t N = 512;                     // FFT boyu
    auto win = hannWindow(N);                 // Kisim 2'deki pencere
    const char* ramp = " .:-=+*#%@";          // yogunluk karakterleri (10 kademe)
    const int nBars = 64;                     // gosterilecek frekans binleri

    for (size_t base = 0; base + N <= x.size(); base += N) {
        std::vector<cf64> buf(N);
        for (size_t n = 0; n < N; ++n) buf[n] = x[base + n] * win[n];
        fft(buf);                             // Kisim 2'deki FFT

        // Ilk nBars bin'i dB'ye cevir ve karaktere esle
        printf("|");
        for (int k = 0; k < nBars; ++k) {
            double mag = std::abs(buf[k]) / N;
            double db = 20.0 * std::log10(mag + 1e-9);
            // -80..0 dB araligini 0..9 karakter indeksine esle
            int level = static_cast<int>((db + 80.0) / 80.0 * 9.0);
            level = std::max(0, std::min(9, level));
            putchar(ramp[level]);
        }
        double tSec = base / fs;
        printf("| t=%.2fs\n", tSec);
    }
    // Ust eksen: her karakter fs/N * (nBars) Hz'e kadar
    printf("0 Hz %*s %.0f Hz\n", nBars - 10, "->", nBars * fs / N);
}

int main() {
    const double fs = 8000.0;
    // Test sinyali: zamanla yukselen bir supurme (chirp) 200 -> 3000 Hz
    std::vector<float> x;
    double phase = 0.0;
    for (int n = 0; n < 8000; ++n) {
        double f = 200.0 + 2800.0 * n / 8000.0;
        phase += 2.0 * M_PI * f / fs;
        x.push_back(0.8f * static_cast<float>(std::sin(phase)));
    }
    spectrumAnalyzer(x, fs);   // ekranda kosegen bir "cizgi" gorursunuz
    return 0;
}
\end{lstlisting}

\ipucu{Bu araç ürettiğiniz her sinyali gözle doğrulamak için harikadır. Bir
chirp'te köşegen çizgi, saf tonda tek dikey sütun, gürültüde homojen bir bulut
görürsünüz. Renkli çıktı için ANSI kaçış kodları (\code{\textbackslash
033[38;5;...m}) ekleyip gerçek bir şelale (waterfall) elde edebilirsiniz.}

% =====================================================================
\gsec{Proje 4 --- Enstrüman Akort Cihazı (Pitch Detection)}
% =====================================================================

Bir gitar veya sesin temel frekansını (pitch) bulup en yakın nota ile kaç
sent (cent) saptığını gösteren bir akort cihazı. Frekansı, FFT yerine
\textbf{otokorelasyonla} buluruz --- gürültülü ve harmonik zengini seslerde daha
güvenilirdir.

\begin{lstlisting}
#include <vector>
#include <cmath>
#include <cstdio>
#include <string>

// Otokorelasyonla temel frekansi kestir. 0 = perde bulunamadi.
double detectPitch(const std::vector<float>& x, double fs,
                   double minHz = 50.0, double maxHz = 1000.0) {
    const size_t minLag = static_cast<size_t>(fs / maxHz);
    const size_t maxLag = static_cast<size_t>(fs / minHz);
    double bestCorr = 0; size_t bestLag = 0;

    // Her olasi periyot (lag) icin otokorelasyonu hesapla
    for (size_t lag = minLag; lag <= maxLag && lag < x.size(); ++lag) {
        double corr = 0;
        for (size_t i = 0; i + lag < x.size(); ++i) corr += x[i] * x[i + lag];
        if (corr > bestCorr) { bestCorr = corr; bestLag = lag; }
    }
    return bestLag ? fs / bestLag : 0.0;
}

// Frekanstan en yakin nota adi + sent sapmasi.
std::string nearestNote(double freq, double* centsOut) {
    static const char* names[] = {"C","C#","D","D#","E","F",
                                  "F#","G","G#","A","A#","B"};
    // A4 = 440 Hz referans; yari-ton = 2^(1/12)
    double midi = 69.0 + 12.0 * std::log2(freq / 440.0);
    int nearest = static_cast<int>(std::lround(midi));
    *centsOut = (midi - nearest) * 100.0;     // sent cinsinden sapma
    int octave = nearest / 12 - 1;
    return std::string(names[(nearest % 12 + 12) % 12]) + std::to_string(octave);
}

int main() {
    const double fs = 44100.0;
    // Test: hafif akortsuz bir "La" (441.5 Hz) uret
    std::vector<float> x;
    for (int n = 0; n < 8192; ++n)
        x.push_back(0.7f * std::sin(2*M_PI*441.5*n/fs)
                  + 0.2f * std::sin(2*M_PI*883.0*n/fs));  // + 2. harmonik

    double f = detectPitch(x, fs);
    double cents;
    std::string note = nearestNote(f, &cents);
    printf("Frekans: %.2f Hz -> Nota: %s (%+.1f sent)\n",
           f, note.c_str(), cents);   // "A4 (+5.9 sent)" -> hafif tiz
    return 0;
}
\end{lstlisting}

\notk[Neden otokorelasyon?]{FFT tepesi temel frekansı verir ama harmonikler
güçlüyse (gitar, insan sesi) yanlış oktava atlayabilir. Otokorelasyon, işaretin
\emph{periyodunu} doğrudan arar ve harmoniklere karşı daha gürbüzdür. Daha da
iyisi \textbf{YIN} ve \textbf{pYIN} algoritmalarıdır (fark fonksiyonu tabanlı);
profesyonel akort uygulamaları bunları kullanır.}

% =====================================================================
\gsec{Projeleri Derleme (CMake)}
% =====================================================================

Dört projeyi de derleyen tek bir CMake dosyası. Her proje bu belgedeki
yardımcıları (\code{writeWav}, \code{fft}, \code{hannWindow}, \code{goertzel},
\code{designLowpass}\dots) tek bir \code{dsp.hpp} başlığında toplayıp
\code{\#include} ettiğinizi varsayar.

\begin{lstlisting}[style=cmakestyle]
cmake_minimum_required(VERSION 3.15)
project(dsp_projeler LANGUAGES CXX)

set(CMAKE_CXX_STANDARD 17)
set(CMAKE_CXX_STANDARD_REQUIRED ON)

# Optimizasyon: DSP kodu SIMD'den cok faydalanir
if(NOT MSVC)
  add_compile_options(-O3 -march=native)
endif()

# Her proje bagimsiz bir calistirilabilir
add_executable(dtmf      proje1_dtmf.cpp)
add_executable(fsk_modem proje2_modem.cpp)
add_executable(spektrum  proje3_spektrum.cpp)
add_executable(akort     proje4_akort.cpp)

# Ortak baslik dizini (dsp.hpp burada)
foreach(t dtmf fsk_modem spektrum akort)
  target_include_directories(${t} PRIVATE ${CMAKE_SOURCE_DIR}/include)
endforeach()
\end{lstlisting}

\begin{lstlisting}[style=shstyle]
# Derle ve calistir
mkdir build && cd build
cmake .. -DCMAKE_BUILD_TYPE=Release
cmake --build .

./dtmf                    # -> Cozulen tus: 5
./fsk_modem               # -> modem.wav uretir, metni geri cozer
./spektrum                # -> terminalde kosegen chirp spektrumu
./akort                   # -> A4 (+5.9 sent)

# Uretilen sesleri dinle / spektrogramla:
sox modem.wav -n spectrogram -o modem.png
\end{lstlisting}

\ipucu{Bu dört projeyi genişletmek için doğal fikirler: DTMF çözücüye gerçek
\code{.wav} girişi ve gürültü reddi; FSK modeme Viterbi FEC ve gerçek ses kartı
G/Ç (PortAudio/RtAudio); spektrum analizörüne ncurses ile canlı kaydırılan
şelale; akort cihazına mikrofon girişi ve grafik iğne göstergesi. Her biri, bu
belgedeki yapı taşlarının üstüne birkaç yüz satırla oturur.}

% #####################################################################
\gpart{9}{Matematiksel Temeller}
% #####################################################################

Bu kısım, belge boyunca kod hâline getirdiğimiz işlemlerin altındaki matematiği
adım adım türetir. Amaç, her algoritmanın \emph{neden} çalıştığını sezgiyle
değil, ispatla görmenizdir. Her sonucu bir önceki koda geri bağlayacağız.

% =====================================================================
\gsec{Karmaşık Üsteller ve Diklik (Orthogonality)}
% =====================================================================

Tüm Fourier teorisi tek bir yapı taşına dayanır: karmaşık üstel
$e^{j\omega t}$. Euler formülü $e^{j\theta} = \cos\theta + j\sin\theta$'dan iki
temel kimlik çıkar:
\[
\cos\theta = \frac{e^{j\theta} + e^{-j\theta}}{2}, \qquad
\sin\theta = \frac{e^{j\theta} - e^{-j\theta}}{2j}
\]
Bu, ``bir kosinüs, biri $+\omega$ diğeri $-\omega$'da dönen iki fazörün
toplamıdır'' demenin cebirsel hâlidir --- gerçek sinüslerin neden hem pozitif
hem negatif frekans içerdiğini (Kısım 4) buradan görürüz.

\gsub{Diklik İlişkisi}

Fourier'in işleyebilmesinin sebebi, farklı frekanstaki üstellerin
\textbf{dik} (birbirine ``bağımsız'') olmasıdır. Sürekli durumda bir $T$
periyodu üzerinde:
\[
\frac{1}{T}\int_0^T e^{j 2\pi m t/T}\, e^{-j 2\pi n t/T}\, dt =
\begin{cases} 1, & m = n \\ 0, & m \neq n \end{cases}
\]
Ayrık durumda ($N$ örnek üzerinde) tam karşılığı:
\[
\sum_{n=0}^{N-1} e^{j 2\pi k n/N}\, e^{-j 2\pi l n/N} =
\begin{cases} N, & k = l \pmod N \\ 0, & \text{diğer} \end{cases}
\]
İkincisi bir \textbf{geometrik seri toplamıdır}: $k \neq l$ için ortak çarpan
$r = e^{j2\pi(k-l)/N} \neq 1$'dir ve $\sum_{n=0}^{N-1} r^n = \frac{r^N - 1}{r - 1}
= 0$ (çünkü $r^N = e^{j2\pi(k-l)} = 1$). İşte DFT'nin \emph{tersinin} neden var
olduğunun kanıtı budur --- her frekans bileşenini diğerlerine karışmadan geri
çekebiliriz.

% =====================================================================
\gsec{Fourier Zinciri: Seriden DFT'ye}
% =====================================================================

Dört Fourier gösterimi vardır; hangisini kullanacağımız işaretin zaman ve
frekansta sürekli mi kesikli mi olduğuna bağlıdır:

\begin{center}
\begin{tabularx}{\textwidth}{@{}>{\bfseries\color{anaMavi}}l >{\raggedright\arraybackslash}X >{\raggedright\arraybackslash}X@{}}
\toprule
\rowcolor{acikMavi}\color{anaMavi}\textbf{Dönüşüm} & \color{anaMavi}\textbf{Zaman} & \color{anaMavi}\textbf{Frekans}\\
\midrule
Fourier Serisi (FS) & sürekli, periyodik & kesikli, sonsuz\\
Fourier Dönüşümü (FT) & sürekli, aperiyodik & sürekli, sonsuz\\
Ayrık-Zaman FT (DTFT) & kesikli, aperiyodik & sürekli, periyodik\\
Ayrık Fourier Dön. (DFT) & kesikli, periyodik & kesikli, periyodik\\
\bottomrule
\end{tabularx}
\end{center}

\gsub{Fourier Serisi ve Dönüşümü}

Periyodu $T$ olan bir işaret, harmoniklerin toplamıdır:
\[
x(t) = \sum_{k=-\infty}^{\infty} c_k\, e^{j 2\pi k t/T}, \qquad
c_k = \frac{1}{T}\int_0^T x(t)\, e^{-j 2\pi k t/T}\, dt
\]
Katsayı formülü doğrudan yukarıdaki diklikten çıkar: her iki tarafı
$e^{-j2\pi l t/T}$ ile çarpıp integre edin, dik olmayan tek terim $k=l$ hayatta
kalır. Periyodu sonsuza götürürsek ($T \to \infty$) toplam integrale, kesikli
$k/T$ frekansları sürekli $f$'e dönüşür ve \textbf{Fourier dönüşümünü} elde
ederiz:
\[
X(f) = \int_{-\infty}^{\infty} x(t)\, e^{-j 2\pi f t}\, dt, \qquad
x(t) = \int_{-\infty}^{\infty} X(f)\, e^{j 2\pi f t}\, df
\]

\gsub{DTFT'den DFT'ye: Kodun Türetimi}

İşareti örnekleyip ($t = nT_s$) kesikli yaparsak DTFT'yi elde ederiz:
\[
X(e^{j\omega}) = \sum_{n=-\infty}^{\infty} x[n]\, e^{-j\omega n}
\]
Burada $\omega = 2\pi f/f_s$ (radyan/örnek, Kısım 1'deki $w$). Bu hâlâ
$\omega$'da süreklidir --- bilgisayarda saklanamaz. Son adım, frekans eksenini de
$N$ noktada örneklemektir: $\omega_k = 2\pi k/N$. Bu bize tam olarak Kısım 2'de
kodladığımız \textbf{DFT}'yi verir:
\[
X[k] = \sum_{n=0}^{N-1} x[n]\, e^{-j 2\pi k n/N}
\]
Yani \code{dft()} fonksiyonu, sürekli Fourier dönüşümünün hem zamanda hem
frekansta örneklenmiş, tam hesaplanabilir hâlidir. Diklik ilişkisi de ters
dönüşümü (\code{idft()}) ve $1/N$ çarpanını doğrular.

\gsub{FFT'nin Böl-ve-Fethet İspatı}

FFT'nin $O(N\log N)$'i, DFT toplamını çift ve tek indekslere ayırmaktan doğar.
$W_N = e^{-j2\pi/N}$ yazarsak:
\[
X[k] = \sum_{n \text{ çift}} x[n] W_N^{kn} + \sum_{n \text{ tek}} x[n] W_N^{kn}
     = \underbrace{\sum_{m=0}^{N/2-1} x[2m] W_{N/2}^{km}}_{E[k]}
     + W_N^k \underbrace{\sum_{m=0}^{N/2-1} x[2m{+}1] W_{N/2}^{km}}_{O[k]}
\]
çünkü $W_N^{2km} = W_{N/2}^{km}$. Ayrıca $W_N^{k+N/2} = -W_N^k$ (simetri)
olduğundan:
\[
X[k] = E[k] + W_N^k O[k], \qquad X[k + N/2] = E[k] - W_N^k O[k]
\]
İşte \code{fftRecursive} içindeki \emph{kelebek} tam olarak budur: iki yarı-boy
DFT hesapla, twiddle çarpanı $W_N^k$ ile birleştir. Her katmanda iş yarıya
inerken $\log_2 N$ katman olması $O(N\log N)$'i verir.

% =====================================================================
\gsec{Örnekleme Teoreminin Türetimi}
% =====================================================================

Nyquist teoremini (Kısım 2) artık ispatlayabiliriz. Örnekleme, işareti bir
\textbf{darbe dizisiyle} (Dirac tarağı) çarpmaktır:
\[
x_s(t) = x(t)\cdot \sum_{n=-\infty}^{\infty}\delta(t - nT_s)
\]
Dirac tarağının Fourier dönüşümü yine bir Dirac tarağıdır (aralık $f_s = 1/T_s$).
\textbf{Çarpımın dönüşümü konvolüsyondur} (bir sonraki bölüm) olduğundan,
örneklenmiş işaretin spektrumu, orijinal spektrumun $f_s$ aralıklarla
\emph{sonsuz kez kopyalanmış} hâlidir:
\[
X_s(f) = f_s \sum_{k=-\infty}^{\infty} X(f - k f_s)
\]
İşte aliasing'in kaynağı budur. Eğer $X(f)$'in bant genişliği $B < f_s/2$ ise
kopyalar üst üste binmez ve orijinali bir alçak geçiren filtreyle kusursuz geri
alabiliriz. Ama $B > f_s/2$ ise komşu kopyalar örtüşür --- yüksek frekanslar
düşüklerin üstüne katlanır (Kısım 2'deki \code{aliasedFrequency}). Kritik sınır
$f_s = 2B$'dir; \textbf{Nyquist oranı} budur.

\notk[İdeal geri çatım]{Örtüşme yoksa, orijinal işaret örneklerden şu
interpolasyonla \emph{tam} olarak geri gelir:
$x(t) = \sum_n x[n]\,\operatorname{sinc}\!\big((t - nT_s)/T_s\big)$.
Bu, kesikli örneklerin neden sürekli işareti eksiksiz temsil edebildiğinin
matematiksel güvencesidir --- sinc çekirdeği ideal alçak geçiren filtrenin dürtü
yanıtıdır, ki Kısım 3'te FIR tasarımında tam bunu kullandık.}

% =====================================================================
\gsec{Konvolüsyon ve Modülasyon Teoremleri}
% =====================================================================

Filtreleme ve modülasyonun tüm matematiği iki eşleştirmeye dayanır. Bunlar
Fourier dönüşümünün en güçlü özellikleridir.

\gsub{Konvolüsyon Teoremi}

Zamanda konvolüsyon, frekansta çarpmadır (ve tersi):
\[
y(t) = (x * h)(t) \;\Longleftrightarrow\; Y(f) = X(f)\, H(f)
\]
İspat, konvolüsyon integralinin dönüşümünü alıp değişken değiştirmektir; $e^{-j2\pi
f\tau}$ çarpanı ayrışır. Bu teoremin iki büyük sonucu var:

\begin{itemize}
\item \textbf{Filtrelemenin anlamı:} Kısım 3'teki FIR/IIR bir konvolüsyondur;
frekans alanında bu, işaretin spektrumunu filtrenin \emph{frekans yanıtı} $H(f)$
ile çarpmaktır. ``Alçak geçiren'' demek, $H(f)$'in düşük frekanslarda $\approx 1$,
yüksekte $\approx 0$ olması demektir.
\item \textbf{Hızlı konvolüsyon:} $O(N^2)$'lik konvolüsyon yerine
$\text{IFFT}\{\text{FFT}\{x\}\cdot\text{FFT}\{h\}\}$ ile $O(N\log N)$'de sonuç ---
uzun FIR'lar için overlap-add yönteminin temeli.
\end{itemize}

\gsub{Modülasyon (Frekans Kaydırma) Teoremi}

Konvolüsyon teoreminin ikizi: zamanda bir üstelle çarpmak, frekansta kaydırmadır:
\[
x(t)\, e^{j 2\pi f_0 t} \;\Longleftrightarrow\; X(f - f_0)
\]
İşte Kısım 4'teki \code{Nco}/\code{frequencyShift}'in ispatı: bir IQ akışını
$e^{-j2\pi f_0 n/f_s}$ ile çarpmak, spektrumunu $-f_0$ kadar öteler ve istediğimiz
kanalı DC'ye getirir. Gerçek bir kosinüsle çarpsaydık ($\cos = \tfrac12(e^{j}
+ e^{-j})$) spektrum \emph{hem} $+f_0$ \emph{hem} $-f_0$'a kopyalanır --- bu
istenmeyen ``görüntü'' (image) tam olarak IQ'nun (tek taraflı spektrum) neden
gerekli olduğunu gösterir.

\gsub{Parseval Teoremi}

Enerji, zaman ve frekans alanında korunur:
\[
\sum_{n=0}^{N-1} |x[n]|^2 = \frac{1}{N}\sum_{k=0}^{N-1} |X[k]|^2
\]
Bu yüzden bir işaretin gücünü spektrumdan ($|X[k]|^2$, yani PSD) hesaplayabiliriz
--- Kısım 2'deki \code{powerSpectrumDb} ve pencere kazancı telafisi bu eşitliğe
dayanır.

% =====================================================================
\gsec{Analitik İşaret ve Anlık Frekans}
% =====================================================================

Son olarak, tüm açı-demodülatörlerinin (FM/PM) altındaki matematiği kesinleştirelim.

\gsub{Anlık Frekansın Tanımı}

Karmaşık bir işareti kutupsal yazalım: $z(t) = a(t)\, e^{j\phi(t)}$. Buradaki
$a(t) = |z(t)|$ \textbf{anlık zarf}, $\phi(t) = \angle z(t)$ \textbf{anlık
fazdır}. \textbf{Anlık frekans}, fazın türevidir:
\[
f_i(t) = \frac{1}{2\pi}\frac{d\phi(t)}{dt}
\]
Bu tanım FM demodülasyonunu doğrudan açıklar: FM'de $\phi(t) = 2\pi k_f\!\int m$,
dolayısıyla $\frac{1}{2\pi}\phi'(t) = k_f\, m(t)$ --- yani fazın türevi tam
olarak mesajdır (bir ölçek dışında).

\gsub{Polar Diskriminatörün Türetimi}

Ayrık zamanda türev, ardışık fazların farkıdır:
$\phi[n] - \phi[n-1]$. Bunu \code{atan2} sarmalanması olmadan hesaplamanın yolu,
Kısım 5'teki eşlenik çarpımıdır. $z[n] = a[n]e^{j\phi[n]}$ için:
\[
z[n]\, z^*[n-1] = a[n]\,a[n-1]\; e^{\,j(\phi[n] - \phi[n-1])}
\]
Genlikler pozitif reel olduğundan bu çarpımın \emph{açısı} tam olarak faz
farkıdır:
\[
\angle\big(z[n]\, z^*[n-1]\big) = \phi[n] - \phi[n-1] \approx 2\pi f_i\, T_s
\]
İşte \code{FmDemod::process}'in matematiksel ispatı budur. \code{std::arg} sonucu
doğal olarak $[-\pi,\pi]$ aralığında verdiği için, faz $\pm\pi$'yi geçse bile
sonuç sarmalanmadan doğru çıkar --- bu yüzden ayrı bir ``unwrap'' adımına gerek
kalmaz.

\gsub{Hilbert Dönüşümü ve Analitik İşaretin Türetimi}

Gerçek bir $x(t)$'nin analitik karşılığı $x_a(t) = x(t) + j\hat{x}(t)$'dir;
burada $\hat{x}$ \textbf{Hilbert dönüşümüdür}. Amaç, negatif frekansları yok edip
tek taraflı bir spektrum elde etmektir. Frekans alanında Hilbert filtresi
$H(f) = -j\,\operatorname{sgn}(f)$'dir; bu, tüm pozitif frekansları $-90^\circ$,
negatifleri $+90^\circ$ döndürür. Bunun sonucunda:
\[
X_a(f) = X(f)\big(1 + \operatorname{sgn}(f)\big) =
\begin{cases} 2X(f), & f > 0 \\ X(0), & f = 0 \\ 0, & f < 0 \end{cases}
\]
İşte Kısım 4'teki \code{analyticSignal} fonksiyonunun ``pozitif frekansları 2 ile
çarp, negatifleri sıfırla'' adımının kaynağı budur. Böylece $|x_a|$ zarfı (AM
demod) ve $\angle x_a$ fazı (FM/PM demod) doğrudan okunabilir hâle gelir --- tüm
demodülatörleri birkaç satıra indiren şey bu matematiktir.

\ipucu{Bütün belgeyi tek cümlede özetlersek: \emph{bir işareti karmaşık bant
temeline taşıdığınızda}, genlik \code{std::abs}, faz \code{std::arg}, frekans ise
ardışık fazların farkıdır. Modülasyon bu üçünden birine bilgi gömer;
demodülasyon onu geri okur. Geri kalan her şey --- örnekleme, FFT, filtreler,
mikser --- bu okumayı temiz ve verimli yapmak içindir.}

\vfill
\begin{center}
{\color{ortaGri}\rule{0.5\linewidth}{0.8pt}}\\[6pt]
{\color{ortaGri}\small Bu rehber C++17 ile, harici bağımlılık olmadan
derlenebilecek örneklerle hazırlanmıştır.\\
Sayısal İşaret İşleme, Modülasyon ve Demodülasyon --- 20 Temmuz 2026}
\end{center}

\end{document}
